\documentclass[12pt]{TA-ITS}
% \documentclass[12pt,twoside]{TA-ITS} % atau book

% SEBELUM ACC (SEBELUM CETAK)
% \pagestyle{draf}

% SETELAH ACC (MAU CETAK)
\pagestyle{TA}

% Singkatan Tugas Akhir
\newcommand{\ta}{penelitian Tugas Akhir}
\newcommand{\Ta}{Penelitian Tugas Akhir}

% Metode & model
\newcommand{\method}{Algoritma Evolusioner dengan Simulasi Agen}
\newcommand{\smethod}{GEEvo-Agent}

% % Singkatan Tugas Akhir
% \newcommand{\ta}{}
% \newcommand{\Ta}{}
% \newcommand{\TA}{}

% % Metode & model
% \newcommand{\method}{}
% \newcommand{\smethod}{}

\begin{document}
    %%%%%%%%%%%%%%%%%%%%%%%%%%%%%%%%%%%%%%%%%%%%%%%%%%%%%%%%%%%%%%%%%%%%%%%%%%%%%%%%
% DATA-DATA UNTUK PERINTAH KHUSUS
% Masukkan data-data sesuai panduan berikut untuk otomatisasi format Tugas Akhir
%%%%%%%%%%%%%%%%%%%%%%%%%%%%%%%%%%%%%%%%%%%%%%%%%%%%%%%%%%%%%%%%%%%%%%%%%%%%%%%%

% Judul Tugas Akhir dalam Bahasa Indonesia & Inggris
\Judul{Optimasi Parameter Sistem Reward pada Ekonomi Gim Menggunakan Algoritma Evolusioner}
\Title{Optimization of Reward System Parameters in Game Economy Using Evolutionary Algorithms}

% Nama mahasiswa yang mengerjakan Tugas Akhir
\Nama{Alridho Tristan Satriawan}

% Nama alias atau nama kecil; jika ada
\NamaKecil{} % Bisa dikosongin, gapapa

% NRP mahasiswa yang mengerjakan Tugas Akhir
\NRP{5002221100}

% Kode mata kuliah Tugas Akhir
\KodeTA{SM234801}

% Asal departemen dalam Bahasa Indonesia & Inggris
\Departemen{Matematika} 
\Department{Mathematics}

% Asal laboratorium, bidang studi atau peminatan; tidak perlu (tahun 2024)
% \Laboratorium{Pembelajaran Mesin dan Big Data}

% Asal fakultas dalam Bahasa Indonesia & Inggris
\Fakultas{Sains dan Analitika Data} % Tanpa kata `fakultas`-nya
\Faculty{Science and Data Analytics} % Tanpa kata `faculty of`-nya`

% Singkatan fakultas dalam Bahasa Indonesia & Inggris
\SingkatanFakultas{FSAD}
\SingkatanFakultasEng{Scientics}

% Tanggal disetujuinya Tugas Akhir
\TanggalDisetujui{30}{Juni}{2026} % {<tanggal>}{<bulan>}{<tahun>}

% Profil kepala departemen
\Kadep{Dr. Didik Khusnul Arif, S.Si., M.Si.} % Nama kepala departemen
\NIPKadep{19730930 199702 1 001} % NIP kepala departemen

% Dosen Pembimbing
\DosenPembimbing{Dr. Drs. Bandung Arry S., MI.Komp.} % Pembimbing pertama
                {} % Pembimbing kedua; abaikan, jika tidak ada
                {} % Pembimbing ketiga; abaikan, jika tidak ada
\NIPPembimbing{19630605 198903 1 003} % NIP Pembimbing pertama
              {} % NIP Pembimbing kedua; abaikan, jika tidak ada
              {} % NIP Pembimbing ketiga; abaikan, jika tidak ada

% Dosen Penguji
\DosenPenguji{Dr. Darmaji, S.Si, MT} % Penguji pertama
             {Soleha, S.Si, M.Si} % Penguji kedua
             {Muhammad Yusuf Hakim Widianto, S.Si., M. Si., M.Sc., Ph.D} % Penguji ketiga
\NIPPenguji{19691015 199412 1 001} % NIP Penguji pertama
           {19830107 200604 2 001} % NIP Penguji kedua
           {199106202024061001} % NIP Penguji ketiga
    \input{tugas-akhir/front-matter/000_front-matter}
    % \input{tugas-akhir/contents/halaman_kosong}
    \chapter*{Abstrak}

\makeatletter
    \vspace{1cm}
    \begin{center}
        \noindent\textbf{\@judul}
    \end{center}

    \vspace{1cm}
    \begin{table}[htbp]
        \begin{tabular}{lcp{0.49\paperwidth}}
            Nama Mahasiswa / NRP & : & \@nama~/ \@NRP\\
            Departemen & : & \@departemen~{\singkatan@fakultas}-ITS\\
            Dosen Pembimbing & : & \pbb@satu
        \end{tabular}
    \end{table}

    \vspace{1cm}
    \noindent
    \textbf{Abstrak}

    % ============================
    % Teks abstrak mulai dari sini
    % ============================
    Permasalahan \textit{reward balancing} dalam desain ekonomi gim merupakan faktor kritis yang memengaruhi keterlibatan dan keberlanjutan pengalaman bermain. Sistem \textit{reward} yang tidak seimbang menyebabkan menurunnya ketertarikan pemain, progres permainan yang tidak wajar, serta instabilitas ekonomi in-gim. Penelitian terdahulu seperti GEEvo \citep{Rupp_2024} telah berhasil mengotomatisasi penyeimbangan ekonomi gim menggunakan algoritma evolusioner, namun belum mempertimbangkan variabilitas perilaku pemain yang bersifat dinamis. Penelitian ini mengembangkan dua kontribusi utama terhadap kerangka GEEvo: (1) simulasi berbasis agen dengan tiga profil perilaku (\textit{aggressive}, \textit{passive}, \textit{random}) untuk merepresentasikan dinamika pemain, dan (2) fungsi fitness obj4 (\textit{tolerance-threshold}) sebagai alternatif yang memberikan gradien sinyal lebih informatif bagi algoritma evolusioner. Evaluasi komparatif pada 194 graf menunjukkan bahwa kombinasi keduanya (Kondisi C) menghasilkan peningkatan \textit{Improved \%} sebesar 87.11\% dan konvergensi pada median 8 generasi, dibandingkan 500 generasi pada kondisi tanpa kontribusi tambahan. Selain itu, dikembangkan pula modul ekstraksi otomatis konfigurasi ekonomi dari \textit{Game Design Document} (GDD) nyata menggunakan pencocokan kata kunci, yang divalidasi pada GDD MMORPG dan VHS \textit{Horror} masing-masing sebanyak 10 graf.

    % Untuk kata kunci
    \katakunci{Algoritma Evolusioner, Simulasi Agen, Reward Balancing, Game Economy, Game Design Document}
\makeatother
    % \input{tugas-akhir/contents/halaman_kosong}
    \chapter*{Abstract}

\makeatletter
    \vspace{1cm}
    \begin{center}
        \noindent\textbf{\@title@}
    \end{center}

    \vspace{1cm}
    \begin{table}[htbp]
        \begin{tabular}{lcp{0.49\paperwidth}}
            Student Name / ID & : & \@nama~/ \@NRP\\
            Department & : & \@department~{\faculty@abbreviation}-ITS\\
            Advisor & : & \pbb@satu
        \end{tabular}
    \end{table}

    \vspace{1cm}
    \noindent
    \textbf{Abstract}
    
    % ====================
    % Abstract starts here
    % ====================
Reward balancing in game economy design is a critical factor affecting player engagement and the sustainability of the gameplay experience. An unbalanced reward system causes declining player interest, unnatural progression, and in-game economic instability. Prior work such as GEEvo \citep{Rupp_2024} has automated game economy balancing using evolutionary algorithms, but has not addressed the variability of player behavior. This research introduces two main contributions to the GEEvo framework: (1) agent-based simulation with three behavioral profiles (aggressive, passive, random) to represent player dynamics, and (2) an obj4 tolerance-threshold fitness function that provides a more informative gradient signal for the evolutionary algorithm. Comparative evaluation on 194 graphs shows that the combined approach (Condition C) achieves an Improved \% of 87.11\% and converges at a median of 8 generations, compared to 500 generations under the baseline. Additionally, an automatic extraction module is developed to derive economy configurations from real Game Design Documents (GDDs) via keyword matching, validated on MMORPG and VHS Horror GDDs with 10 graphs each.


    % For keywords
     \keywords{Evolutionary Algorithm, Agent-based Simulation, Reward Balancing, Game Economy, Game Design Document}
\makeatother
    % \input{tugas-akhir/contents/halaman_kosong}
    \chapter*{Kata Pengantar}
Puji syukur penulis panjatkan kepada Tuhan Yang Maha Esa karena atas berkat dan anugerah-Nya sehingga penulis dapat menyelesaikan Tugas Akhir yang berjudul:
\begin{center}
    \makeatletter
    ``\noindent\MakeUppercase{\@judul}''
    \makeatother
\end{center}

\noindent \Ta~ini disusun guna memenuhi salah satu syarat kelulusan pada jenjang Sarjana di Departemen Matematika, Fakultas Sains dan Analitika Data, Institut Teknologi Sepuluh Nopember. \Ta~ini tentu tidak akan selesai tanpa dukungan, kerja sama, dan bantuan dari banyak pihak. Oleh karena itu, penulis ingin menyampaikan rasa terima kasih yang mendalam kepada:
%
\begin{enumerate}
    \item Kedua orang tua penulis yang senantiasa memberikan doa, dukungan, dan dorongan selama perkuliahan.
    \item Bapak Dr. Didik Khusnul Arif, S.Si., M.Si., selaku Kepala Departemen Matematika FSAD ITS atas dukungan sarana prasarana yang diberikan kepada penulis.
    \item Bapak Dr. Drs. Bandung Arry S., MI.Komp., selaku dosen pembimbing yang selalu memberikan motivasi, arahan, dukungan moral dan fisik selama penyusunan \ta.
    \item Bapak/Ibu dosen pengajar yang tidak dapat disebutkan satu per satu atas ilmu yang diberikan, serta seluruh staf, tenaga kependidikan, dan keluarga besar Departemen Matematika, Fakultas Sains dan Analitika Data, ITS, atas bantuan dan dukungan selama perkuliahan berlangsung.
    \item Teman-teman seperjuangan yang telah senantiasa menemani serta memberikan dukungan dan bantuan selama perkuliahan.
\end{enumerate}

Penulis menyadari bahwa \ta~masih jauh dari kata sempurna. Oleh karena itu, segala bentuk masukan maupun saran yang membangun dari para pembaca sangat diharapkan guna untuk pengembangan penulisan kedepannya. Akhir kata, sekiranya \ta~ dapat memberikan manfaat bagi semua pihak yang berkepentingan.

\makeatletter
\begin{flushright}
    Surabaya, \ \@tanggal\ \@bulan\ \@tahun

    \vspace{2cm}
    \@nama
\end{flushright}
\makeatother
    \newpage
    % \input{tugas-akhir/contents/halaman_kosong}
    \listoffigures
    % \input{tugas-akhir/contents/halaman_kosong}
    \listoftables
    % \input{tugas-akhir/contents/halaman_kosong}
    \listofalgorithms
    % \input{tugas-akhir/contents/halaman_kosong}
    \chapter*{Daftar Simbol}
\begin{center}
\small % Kurangi ukuran font
\renewcommand{\arraystretch}{1.05} % Kurangi jarak baris dari 1.2
\setlength{\tabcolsep}{4pt} % Kurangi padding dari 6pt
\begin{tabular}{p{2.5cm}p{0.2cm}p{10cm}} % Kurangi lebar kolom pertama
$G = (V, E)$ & : & Graf berarah yang merepresentasikan ekonomi gim dengan himpunan vertex $V$ dan himpunan edge $E$ \\
$V$ & : & Himpunan vertex (node) dalam graf ekonomi gim \\
$E$ & : & Himpunan edge (sisi) berarah dalam graf \\
$v_i$ & : & Vertex ke-$i$ dalam graf, merepresentasikan komponen ekonomi \\
$n$ & : & Jumlah node dalam graf ekonomi gim \\
$w_{ij}$ & : & Bobot aliran sumber daya dari node $i$ ke node $j$ \\
$W$ & : & Matriks bobot aliran antar node: $W = [w_{ij}]_{n \times n}$ \\
$\text{Outflow}(v_i)$ & : & Jumlah total aliran keluar dari node $v_i$: $\sum_{j=1}^{n} w_{ij}$ \\
$\text{Inflow}(v_i)$ & : & Jumlah total aliran masuk ke node $v_i$: $\sum_{j=1}^{n} w_{ji}$ \\
$f(G)$ & : & Fungsi objektif untuk optimasi keseimbangan graf \\
$P_0$ & : & Populasi awal dalam algoritma evolusioner \\
$P_t$ & : & Populasi pada generasi ke-$t$ \\
$P_{t+1}$ & : & Populasi pada generasi ke-$(t+1)$ hasil evolusi \\
$x_i$ & : & Individu (solusi kandidat) ke-$i$ dalam populasi \\
$x^*$ & : & Solusi optimal hasil algoritma evolusioner \\
$f_{\text{gen}}(G)$ & : & Fungsi fitness pada tahap \textit{generator} untuk mengevaluasi validitas struktur graf ekonomi gim \\
$\mathrm{validate}(v)$ & : & Fungsi bantu yang menghitung jumlah batasan yang tidak terpenuhi oleh node $v$ \\
$f_1(s^{p_j}_t, x)$ & : & Fungsi fitness pada tahap \textit{balancer} untuk mengukur kedekatan nilai sumber daya terhadap target absolut \\
$s^{p_j}_t$ & : & Jumlah sumber daya dalam pool $p_j$ pada waktu ke-$t$ hasil simulasi ekonomi gim \\
$p_j$ & : & Pool sumber daya ke-$j$ dalam graf ekonomi gim \\
$x$ & : & Nilai target absolut yang ingin dicapai oleh sumber daya (misal jumlah \textit{gold}) \\
$\alpha$ & : & Parameter ambang batas (\textit{threshold}) untuk toleransi rata-rata hasil simulasi pada fungsi fitness \\
$m$ & : & Jumlah simulasi yang dijalankan untuk mengurangi pengaruh keacakan (\textit{randomness}) \\
$i$ & : & Indeks simulasi ke-$i$ dalam perhitungan rata-rata fitness \\
$\mathrm{prop}(s^{p_j}_t, x)$ & : & Fungsi proporsional untuk mengukur kedekatan nilai sumber daya terhadap target \\
$N$ & : & Ukuran populasi dalam algoritma evolusioner \\
$f(\cdot)$ & : & Fungsi fitness untuk evaluasi kualitas individu \\
$p_c$ & : & Probabilitas crossover dalam algoritma evolusioner \\
$p_m$ & : & Probabilitas mutasi dalam algoritma evolusioner \\
$x_i'$ & : & Individu hasil mutasi dari $x_i$ \\
$\epsilon$ & : & Variabel acak dari distribusi normal untuk mutasi \\
$\sigma^2$ & : & Variansi untuk operasi mutasi \\
$\mathcal{N}(0, \sigma^2)$ & : & Distribusi normal dengan mean 0 dan variansi $\sigma^2$ \\
\end{tabular}
\end{center}
% \newpage
% \begin{center}
% \renewcommand{\arraystretch}{1.2}
% \setlength{\tabcolsep}{6pt}
% \begin{tabular}{p{3cm}p{0.3cm}p{11cm}}

% % $O(\cdot)$ & : & Notasi kompleksitas waktu asimtotik \\
% \end{tabular}
% \end{center}



% \begin{table}[htbp]
%     \begin{tabular}{lcl}
%         $\mu$ & : & Nilai rata-rata\\
%         $\sigma$ & : & Nilai standar deviasi\\
%     \end{tabular}
% \end{table}
% \clearpage
% \phantomsection
% \addcontentsline{toc}{chapter}{Daftar Simbol}

% % \newgeometry{left=2.2cm,right=2.2cm,top=3cm,bottom=3cm}

% \vspace*{-1cm} % naikkan posisi jika terlalu ke bawah
% \begin{center}
% \Large\bfseries DAFTAR SIMBOL
% \end{center}

% \vspace{0.8cm}

% \small
% \renewcommand{\arraystretch}{1.2}
% \setlength{\tabcolsep}{6pt}

% \begin{center}
% \begin{tabular}{p{3cm} p{11.5cm}}
% $G = (V,E,w)$ & Graf berbobot dengan himpunan simpul $V$, himpunan sisi $E$, dan fungsi bobot $w$ \\
% $N, A$ & Himpunan simpul (nodes) dan busur (arcs) pada graf \\
% $s, g$ & Simpul awal (\textit{start}) dan simpul tujuan (\textit{goal}) \\
% $\Gamma$ & Himpunan simpul tujuan \\
% $c(u,v)$ & Biaya (cost) dari simpul $u$ ke simpul $v$ \\
% $\vec{c}(u,v)$ & Vektor biaya berdimensi jamak (\textit{multi-objective cost vector}) \\
% $\vec{g}(n)$ & Vektor biaya aktual dari simpul awal ke simpul $n$ \\
% $g(n)$ & Biaya aktual kumulatif dari simpul awal ke simpul $n$ \\
% $h(n)$ & Nilai heuristik (estimasi biaya dari simpul $n$ ke tujuan) \\
% $H(n)$ & Himpunan vektor heuristik pada simpul $n$ dalam MOA* \\
% $f(n)$ & Fungsi evaluasi total: $f(n) = g(n) + h(n)$ \\
% $F(P_{sn})$ & Himpunan nilai evaluasi tidak terdominasi dari jalur $P_{sn}$ \\
% $\vec{f}$ & Vektor fungsi evaluasi dalam ruang multiobjektif \\
% $\prec$ & Relasi dominasi Pareto (lebih baik pada setidaknya satu dimensi dan tidak lebih buruk pada lainnya) \\
% $\text{ND}(X)$ & Himpunan vektor tidak terdominasi dari himpunan $X$ \\
% $\mathcal{S}^*$ & Himpunan solusi Pareto-optimal \\
% $G_{op}(n)$ & Label aktif pada simpul $n$ yang belum diekspansi (OPEN set) \\
% $G_{cl}(n)$ & Label yang telah dievaluasi (CLOSED set) \\
% $\mathcal{G}$ & Himpunan solusi Pareto sementara dari setiap \textit{thread} dalam PMOA* \\
% $\mathbf{UB}$ & Vektor batas atas global untuk setiap dimensi biaya \\
% $d$ & Jumlah objektif atau dimensi biaya dalam MOA* / PMOA* \\
% $\mathbf{c}_i$ & Urutan vektor biaya yang dirotasi siklik untuk \textit{thread} ke-$i$ \\
% $\lceil \mathbf{c} \rceil_k$ & Operator \textit{truncation} untuk reduksi dimensi hingga $k$ kali \\
% $\text{OPEN}$ & Himpunan simpul yang ditemukan tetapi belum dievaluasi \\
% $\text{CLOSED}$ & Himpunan simpul yang telah dievaluasi \\
% $\text{Pareto set}$ & Himpunan solusi tidak terdominasi (\textit{Pareto front}) \\
% $\text{BBST}$ & Struktur data \textit{Balanced Binary Search Tree} (digunakan pada EMOA*) \\
% PMOA* & \textit{Parallel Multi-Objective A* Search} \\
% MOA* & \textit{Multi-Objective A* Search} \\
% NAMOA* & \textit{New Approach to Multi-Objective A*} \\
% BOA* & \textit{Bi-Objective A*} \\
% EMOA* & \textit{Enhanced Multi-Objective A*} \\
% LTMOA* & \textit{Lazy Truncated Multi-Objective A*} \\
% NWMOA* & \textit{Non-Weighted Multi-Objective A*} \\
% $O(\cdot)$ & Notasi kompleksitas waktu asimtotik \\
% \end{tabular}
% \end{center}

% \restoregeometry


% % \usepackage{tabularx}
% \clearpage  
% \chapter*{Daftar Simbol}
% \begin{center}
% \begin{tabularx}{\textwidth}{l c X}
% $G = (V, E, w)$ & : & Graf berbobot dengan himpunan simpul $V$, himpunan sisi $E$, dan fungsi bobot $w$ \\[3pt]
% $N, A$ & : & Himpunan simpul (\textit{nodes}) dan busur (\textit{arcs}) pada graf \\[3pt]
% $s, g$ & : & Simpul awal (\textit{start}) dan simpul tujuan (\textit{goal}) \\[3pt]
% $\Gamma$ & : & Himpunan simpul tujuan \\[3pt]
% $c(u,v)$ & : & Biaya (\textit{cost}) dari simpul $u$ ke simpul $v$ \\[3pt]
% $\vec{c}(u,v)$ & : & Vektor biaya multidimensi (\textit{multi-objective cost vector}) \\[3pt]
% $\vec{g}(n)$ & : & Vektor biaya aktual dari simpul awal ke simpul $n$ \\[3pt]
% $g(n)$ & : & Biaya aktual kumulatif dari simpul awal ke simpul $n$ \\[3pt]
% $h(n)$ & : & Nilai heuristik (estimasi biaya dari simpul $n$ ke simpul tujuan) \\[3pt]
% $H(n)$ & : & Himpunan vektor heuristik pada simpul $n$ dalam MOA* \\[3pt]
% $f(n)$ & : & Fungsi evaluasi total: $f(n) = g(n) + h(n)$ \\[3pt]
% $F(P_{sn})$ & : & Himpunan nilai evaluasi tidak terdominasi dari jalur $P_{sn}$ \\[3pt]
% $\vec{f}$ & : & Vektor fungsi evaluasi dalam ruang multiobjektif \\[3pt]
% $\prec$ & : & Relasi dominasi Pareto (lebih baik pada satu dimensi dan tidak lebih buruk pada lainnya) \\[3pt]
% $\text{ND}(X)$ & : & Himpunan vektor tidak terdominasi dari himpunan $X$ \\[3pt]
% $\mathcal{S}^*$ & : & Himpunan solusi Pareto-optimal \\[3pt]
% $G_{op}(n)$ & : & Label aktif pada simpul $n$ yang belum diekspansi (\textit{OPEN set}) \\[3pt]
% $G_{cl}(n)$ & : & Label yang telah dievaluasi (\textit{CLOSED set}) \\[3pt]
% $\mathcal{G}$ & : & Himpunan solusi Pareto sementara dari setiap \textit{thread} dalam PMOA* \\[3pt]
% $\mathbf{UB}$ & : & Vektor batas atas global untuk setiap dimensi biaya \\[3pt]
% $d$ & : & Jumlah objektif atau dimensi biaya dalam MOA*/PMOA* \\[3pt]
% $\mathbf{c}_i$ & : & Urutan vektor biaya yang dirotasi siklik untuk \textit{thread} ke-$i$ \\[3pt]
% $\lceil \mathbf{c} \rceil_k$ & : & Operator \textit{truncation} untuk reduksi dimensi hingga $k$ kali \\[3pt]
% $\text{OPEN}$ & : & Himpunan simpul yang ditemukan tetapi belum dievaluasi \\[3pt]
% $\text{CLOSED}$ & : & Himpunan simpul yang telah dievaluasi \\[3pt]
% $\text{Pareto set}$ & : & Himpunan solusi tidak terdominasi (\textit{Pareto front}) \\[3pt]
% $\text{BBST}$ & : & Struktur data \textit{Balanced Binary Search Tree} (digunakan pada PMOA*) \\[3pt]
% PMOA* & : & \textit{Parallel Multi-Objective A* Search} \\[3pt]
% MOA* & : & \textit{Multi-Objective A* Search} \\[3pt]
% NAMOA* & : & \textit{New Approach to Multi-Objective A*} \\[3pt]
% BOA* & : & \textit{Bi-Objective A*} \\[3pt]
% EMOA* & : & \textit{Enhanced Multi-Objective A*} \\[3pt]
% LTMOA* & : & \textit{Lazy Truncated Multi-Objective A*} \\[3pt]
% NWMOA* & : & \textit{Non-Weighted Multi-Objective A*} \\[3pt]
% $O(\cdot)$ & : & Notasi kompleksitas waktu asimtotik \\
% \end{tabularx}
% \end{center}
% % \addcontentsline{toc}{chapter}{Daftar Simbol}




    % \input{tugas-akhir/contents/halaman_kosong}
    \newpage
    \tableofcontents
    \newpage
    \setstretch{1.2}
    \pagenumbering{arabic}
    \setcounter{page}{1}
    \chapter{Pendahuluan}\label{ch:introduction}

% Latar Belakang
\section{Latar Belakang Masalah}\label{sec:background}

Dalam dunia industri gim, desain gim menjadi faktor kunci dalam menentukan kesuksesan dan keberlanjutan sebuah gim. Hal ini ditunjukkan oleh survei global terhadap pengembang gim yang dilakukan oleh Statista, tercatat bahwa 58 persen responden menilai kemajuan desain game sebagai faktor penting bagi pertumbuhan industri gim di masa depan \citep{IGDA2019}. Temuan ini menegaskan bahwa keputusan desain berpean besar dalam membentuk kualitas sebuah gim. Salah satu komponen desain gim yang memiliki pengaruh langsung terhadap pengalaman pemain adalah sistem \textit{reward}, yang merupakan subset dari desain ekonomi gim. Sistem ini berperan dalam mengatur hubungan antar usaha yang dilakukan pemain dengan hasil yang diperoleh. Ketidakseimbangan pada sistem \textit{reward} dapat menyebabkan pemain kehilangan minat, progres permainan menjadi terlalu cepat atau lambat, serta menurunkan stabilitas ekonomi \textit{in-game} \citep{Holmberg2021}. Pengaturan sistem \textit{reward} secara manual akan semakin sulit dilakukan seiring dengan meningkatnya kompleksitas desain game dan perilaku pemain. Oleh karena itu, diperlukan pengoptimalan sistem \textit{reward} yang mampu beradaptasi agar tetap seimbang, adil, dan meningkatkan peluang keberlanjutan pengalaman bermain \citep{Rupp_2024}.

Beberapa penelitian terdahulu telah mengkaji berbagai pendekatan untuk menyeimbangkan ekonomi gim dan sistem \textit{reward}. Pada tahap awal, sistem ekonomi gim dirancang, disesuaikan, serta diuji secara manual oleh desainer melalui pengaturan langsung \citep{schreiber2021game}. Pendekatan intuitif ini sulit untuk mempertahankan keseimbangan ekosistem gim yang kompleks dan dinamis. Seiring perkembangan penelitian, muncul pendekatan otomatis berbasis \textit{procedural content generation} (PCG) seperti yang dilakukan oleh \citet{sorochan2022generatingrealtimestrategygame} dalam studi mereka \textit{Generating Real-Time Strategy Game Units Using Search-Based Procedural Content Generation and Monte Carlo Tree Search}, yang menggabungkan pencarian berbasis Monte Carlo untuk menghasilkan dan menyeimbangkan unit gim strategi secara adaptif. Pendekatan tersebut menandai dengan dimulainya pergeseran menuju otomasi desain dengan pemanfaatan algoritma pencarian cerdas. Selanjutnya terdapat penelitian dari \citet{Rupp_2024} dalam judul \textit{GEEvo : Game Economy Generation and Balancing with Evolutionary Algorithms} yang memperluas dengan menerapkan algoritma evolusioner untuk menghasilkan serta menyeimbangkan ekonomi gim secara menyeluruh. Di sisi lain, penelitian \textit{Optimizing Hearthstone Agents using an Evolutionary Algorithm} oleh \cite{garcia2020optimizing} menunjukkan efektivitas algoritma evolusioner dalam mengoptimasi perilaku agen gim melalui proses \textit{co-evolutionary training}, yang relevan dalam konteks optimasi sistem \textit{reward} adaptif berbasis interaksi pemain. Secara keseluruhan, perkembangan penelitian ini menunjukkan pergeseran dari metode manual menuju otomatis yang mampu menyeimbangkan dan menyesuaikan ekonomi gim secara dinamis.  

Meskipun berbagai penelitian telah berupaya menyeimbangkan sistem \textit{reward} dan ekonomi gim menggunakan pendekatan algoritma, sebagian besar masih memiliki keterbatasan dalam adaptivitas terhadap perubahan perilaku pemain yang dinamis. Pendekatan berbasis \textit{procedural content generation} seperti yang dikembangkan oleh Sorochan dan Guzdial (2022) berfokus pada generasi konten gim, bukan keseimbangan sistem \textit{reward} yang berinteraksi langsung dengan perilaku pemain. Di sisi lain, penelitian seperti GEEvo \citep{Rupp_2024} telah berhasil mengotomatisasi sebagian proses perancangan ekonomi gim menggunakan algoritma evolusioner, namun belum mempertimbangkan variabilitas preferensi dan strategi pemain dalam skala mikro. Pendekatan lain seperti optimasi agen gim menggunakan \textit{co-evolutionary algorithms} \citep{garcia2020optimizing} memperlihatkan kemampuan adaptif terhadap dinamika strategi, tetapi belum diintegrasikan dengan keseluruhan sistem ekonomi gim. Selain itu, penelitian terdahulu umumnya masih menggunakan model statis yang tidak memperhitungkan perubahan nilai \textit{reward} seiring evolusi ekonomi dan tingkat keterlibatan pemain \citep{Bailey2024}. Oleh karena itu, diperlukan pendekatan baru yang mampu memadukan adaptivitas perilaku pemain dengan optimasi algoritmik secara evolusioner, sehingga sistem \textit{reward} dapat menyesuaikan diri terhadap perubahan parameter ekonomi gim dan menjaga keseimbangannya secara berkelanjutan.

Berdasarkan keterbatasan dari penelitian sebelumnya, maka diperlukan pendekatan yang tidak hanya berfokus pada ekonomi gim secara makro, tetapi juga mampu beradaptasi terhadap variasi perilaku pemain. Maka dari itu, penelitian ini berfokus untuk mengembangkan model optimasi sistem \textit{reward} berbasis algoritma evolusioner dengan kombinasi simulasi agen untuk mempresentasikan interaksi dan dinamika pemain di dalam gim. Tantangan dalam pendekatan ini terletak pada bagaimana model dapat menjaga sistem \textit{reward} tetap stabil meskipun interaksi multi-agen mengalami kompleksitas yang bersifat stokastik. Selain itu fungsi fitness juga diperlukan untuk menyeimbangkan \textit{fairness, engagement,} serta \textit{sustainability} dalam ekosistem gim. Dengan demikian, penelitian ini berupaya memfokuskan optimasi sistem ekonomi yang bersifat makro menjadi model adaptif berbasis perilaku pemain yang lebih representatif terhadap dinamika ekonomi gim. penelitian ini diharapkan dapat menjadi langkah awal menuju model ekonomi gim adaptif yang tidak hanya menyeimbangkan nilai numerik, tetapi juga mengintegrasikan aspek psikologis dan perilaku pemain dalam proses optimasinya.


% Beberapa penelitian sebelumnya telah dikembangkan untuk menangkap gerakan tubuh manusia. Metode berbasis video dan kamera (misalnya motion capture optik) memerlukan konfigurasi dan mudah mengalami masalah okklusi, sehingga kurang cocok untuk penggunaan luas. merupakan pendekatan pertama yang mengombinasikan neural kinematics estimator dan physics-aware motion optimizer untuk melacak gerakan tubuh hanya dengan enam sensor IMU. Metode ini mampu merekonstruksi pose 3D, torsi sendi, serta gaya reaksi tanah secara real-time dengan akurasi, stabilitas temporal, dan konsistensi fisik. Selanjutnya, Jiang et al. (2022) mengusulkan Transformer Inertial Poser (TIP) yang memanfaatkan teknologi model transformer untuk meningkatkan konsistensi temporal prediksi gerak dan mengevaluasi drift pada posisi global maupun sendi selama rekonstruksi. TIP juga menambahkan modul pemodelan lingkungan untuk meningkatkan realisasi gerak pada berbagai keadaan. Terakhir, IMUPoser yang dikembangkan oleh Mollyn et al. (2023) menekankan penggunaan sensor IMU dalam perangkat sehari-hari seperti ponsel, jam tangan pintar, dan earbud, sehingga memungkinkan pelacakan pose tubuh tanpa memerlukan perangkat khusus tambahan yang mahal. IMUPoser mampu menerima beragam kombinasi input sensor yang tersedia dan tetap menghasilkan estimasi pose tubuh secara real-time.

% Meskipun perkembangan rekonstruksi gerak manusia mengalami peningkatan, terdapat beberapa keterbatasan pada metode sebelumnya yang mendasari research gap saat ini. Metode berbasis IMU yang dilakukan oleh penelitian terdahulu mengharuskan konfigurasi sensor yang tetap sehingga kurang fleksibel terhadap perubahan konfigurasi (Van Wouwe et al., 2024). Jika terdapat perubahan konfigurasi salah satu sensor, akurasi dapat menurun drastis dan model perlu dilatih ulang. Selain itu, model regresif konvensional tidak mampu mengisi kekosongan data ketika terjadi hilangnya sinyal sensor secara real-time, akibatnya gerakan yang direkonstruksi dapat terputus atau tidak alami saat data input tidak lengkap. Beberapa pendekatan juga masih mengalami drift pada estimasi posisi jangka panjang serta kesulitan menangani keragaman jenis gerakan dan lingkungan (Jiang et al., 2022). IMUPoser, meskipun lebih fleksibel tetapi harus menghadapi tantangan data sensor yang bising dan jumlah atau posisi sensor yang dinamis dari waktu ke waktu (Mollyn et al., 2023). Keterbatasan-keterbatasan ini menegaskan perlunya sebuah sistem rekonstruksi gerak yang lebih adaptif terhadap konfigurasi sensor, mampu mengatasi kehilangan data, rendah latensi, dan tetap menjaga gerakan sesungguhnya.

% Untuk mengatasi gap tersebut, penelitian ini berfokus pada penggunaan dan pengembangan model DiffusionPoser (Van Wouwe et al., 2024) sebagai inti sistem rekonstruksi gerakan untuk dunia virtual. DiffusionPoser merupakan model generatif berbasis diffusion yang mampu merekonstruksi gerakan tubuh manusia secara real-time dari kombinasi sensor tanpa memerlukan pelatihan ulang model untuk tiap konfigurasi. Model ini memberikan kelebihan berupa fleksibilitas dalam menentukan jumlah dan penempatan sensor yang dipakai sekaligus memastikan keluaran gerak yang dihasilkan tetap realistis meskipun terdapat kebebasan gerak yang tidak teramati langsung oleh sensor. Sifat generatif difusi dalam DiffusionPoser berperan sebagai acuan yang dapat mengisi gerakan yang hilang atau tidak tertangkap sensor, sehingga avatar virtual dapat bergerak lancar dan sesuai dengan kenyataan aslinya. Beberapa tantangan teknis perlu diperhatikan agar model terintegrasi ke dalam ekosistem dunia virtual, antara lain menjaga latensi tetap rendah, kalibrasi dan sinkronisasi berbagai jenis sensor wearable yang mungkin digunakan, serta memastikan kesesuaian gerakan hasil rekonstruksi dengan gerakan asli pengguna tanpa distorsi yang mengganggu pengalaman imersif.

% Penelitian tugas akhir ini mengusulkan pengembangan dan integrasi model generatif difusi berbasis DiffusionPoser ke dalam ekosistem dunia virtual interaktif. Solusi yang diusulkan diharapkan mampu menghadirkan sistem pelacakan gerakan tubuh yang lebih fleksibel, adaptif, dan realistis, sehingga memberikan pengalaman imersif. Dengan kata lain, karya ini bertujuan menjembatani research gap yang ada dengan menghadirkan platform rekonstruksi gerak real-time untuk dunia virtual yang mengombinasikan fleksibilitas konfigurasi sensor dan realisme gerakan hasil model difusi generatif terkini.



% Rumusan Masalah
\section{Rumusan Masalah}
\label{sec:problem}

Rumusan masalah dalam penelitian Tugas Akhir ini dinyatakan sebagai berikut:

\begin{enumerate}

     \item Bagaimana merancang model sistem \textit{reward} dalam ekonomi gim sehingga mampu merepresentasikan aliran sumber daya?

     \item Bagaimana mengembangkan simulasi berbasis agen untuk merepresentasikan perilaku pemain terhadap sistem \textit{reward} dalam ekonomi gim?

    \item Bagaimana mengembangkan integrasi algoritma evolusioner untuk mengoptimasi parameter sistem \textit{reward} berdasarkan hasil simulasi agen?
    
    
\end{enumerate}

% Batasan Masalah
\section{Batasan Masalah}
\label{sec:limitation}

Penelitian ini dibatasi pada sistem \textit{reward} dalam ekonomi gim berbasis progresi sumber daya, seperti gim \textit{role-playing game} (RPG) atau \textit{simulation game}, di mana pemain memperoleh reward berupa mata uang virtual, pengalaman (XP), atau item melalui aktivitas tertentu, serta tidak berfokus pada implementasi pada gim komersial tertentu, melainkan menggunakan model ekonomi gim abstrak sebagai studi kasus.

% Tujuan Masalah
\section{Tujuan Penelitian}
\label{sec:purpose}

Adapun tujuan dari penelitian Tugas Akhir ini sebagai berikut : 
%     Tujuan Penelitian Tugas Akhir
     \begin{enumerate}

         \item Merancang model sistem \textit{reward} dalam ekonomi gim yang mampu merepresentasikan aliran sumber daya

         \item Mengembangkan simulasi berbasis agen untuk merepresentasikan perilaku pemain terhadap sistem \textit{reward} dalam ekonomi gim

         \item Mengembangkan integrasi algoritma evolusioner untuk mengoptimasi parameter sistem \textit{reward} berdasarkan hasil simulasi agen
        
     \end{enumerate}

% Manfaat Penelitian
\section{Manfaat}
\label{sec:benefit}

Manfaat dari penelitian Tugas Akhir ini dijabarkan sebagai berikut:

\begin{enumerate}
    \item Menyediakan kerangka eksperimental yang dapat mendukung pengembangan sistem otomatis dalam penyeimbangan ekonomi gim, sehingga mengurangi ketergantungan pada pengujian manual oleh desainer gim.
    \item Memberikan pendekatan komputasional berbasis algoritma evolusioner untuk mengoptimasi sistem \textit{reward} dalam ekonomi gim, yang dapat menjadi referensi bagi penelitian lanjutan
\end{enumerate}

    
    \chapter{Tinjauan Pustaka}\label{ch:literature}

\section{Penelitian Terdahulu}\label{sec:past research}

Keseimbangan sistem reward dalam ekonomi gim merupakan faktor penting dalam menjaga kualitas pengalaman bermain. Sistem reward yang tidak seimbang dapat menyebabkan pemain kehilangan motivasi, mengalami kebosanan, atau merasa frustrasi, sehingga berdampak pada keberlanjutan permainan \citep{Holmberg2021}. Desain ekonomi gim, termasuk mekanisme pemberian hadiah, progresi level, dan distribusi sumber daya virtual, berpengaruh langsung terhadap dinamika permainan dan kepuasan pemain.

Kompleksitas ekonomi gim modern menyebabkan proses penyeimbangan menjadi sulit dilakukan secara manual. Perubahan kecil pada parameter ekonomi sering kali menimbulkan dampak signifikan dan tidak terduga terhadap keseluruhan sistem \citep{Volz2016}. Oleh karena itu, berbagai penelitian mengusulkan penggunaan algoritma evolusioner untuk mengotomatisasi proses optimasi dan \textit{balancing} sistem \textit{reward}.

\cite{Volz2016} mendemonstrasikan kelayakan \textit{automatic game balancing} melalui optimasi parameter kartu pada permainan kartu menggunakan algoritma evolusioner berbasis simulasi. Pendekatan ini mampu menghasilkan konfigurasi parameter yang seimbang tanpa memerlukan penalaan manual yang ekstensif. Selanjutnya, de Mesentier Silva et al. \cite{deMesentierSilva2019} menerapkan algoritma evolusioner untuk menyeimbangkan meta permainan Hearthstone dengan mengevolusi atribut kartu, sehingga tidak ada strategi yang terlalu dominan.

\cite{garcia2020optimizing} mengembangkan pendekatan optimasi agen gim Hearthstone menggunakan algoritma evolusioner. Penelitian mereka menunjukkan efektivitas algoritma evolusioner dalam mengoptimasi perilaku agen gim melalui proses \textit{co-evolutionary training}. Pendekatan ini relevan dalam konteks optimasi sistem reward adaptif berbasis interaksi pemain, meskipun masih berfokus pada strategi permainan individual dan belum terintegrasi dengan keseluruhan sistem ekonomi gim.

Pada genre lain, \cite{Gravina2016} memperkenalkan \textit{Constraint Surprise Search} untuk menghasilkan senjata yang beragam namun tetap seimbang dalam gim \textit{first person shooter}. Sementara itu, \cite{Morosan2017} menunjukkan bahwa algoritma evolusioner dapat digunakan untuk menyeimbangkan unit dalam gim \textit{real time strategy} dengan menyesuaikan parameter unit secara otomatis.

\cite{sorochan2022generatingrealtimestrategygame} mengembangkan metode berbasis \textit{procedural content generation} (PCG) yang menggabungkan pencarian berbasis Monte Carlo untuk menghasilkan dan menyeimbangkan unit gim strategi secara adaptif. Pendekatan ini menandai pergeseran menuju otomasi desain dengan pemanfaatan algoritma pencarian cerdas. Namun, pendekatan tersebut masih berfokus pada generasi konten gim dan belum mempertimbangkan keseimbangan sistem reward yang berinteraksi langsung dengan perilaku pemain secara komprehensif.

% Sebagian besar penelitian awal berfokus pada komponen permainan secara individual. Namun, \cite{Rogers2023} memperluas pendekatan ini dengan mengoptimasi ekonomi gim secara keseluruhan. Mereka merepresentasikan ekonomi gim sebagai graf alur sumber daya dan menggunakan algoritma evolusioner untuk menargetkan tingkat kompleksitas ekonomi tertentu. Pendekatan ini memungkinkan kontrol yang lebih terstruktur terhadap sistem reward dan ekonomi gim.

Kerangka kerja \textit{Game Economy Evolution} (GEEvo) yang dikembangkan oleh Rupp \& Eckert \cite{Rupp_2024} selanjutnya membagi proses perancangan ekonomi gim menjadi tahap generasi struktur dan tahap optimasi parameter. Algoritma evolusioner digunakan untuk menyesuaikan parameter ekonomi berdasarkan tujuan \textit{balancing} yang ditetapkan, sehingga proses iterasi desain dapat dilakukan secara sistematis dan fleksibel. Meskipun GEEvo telah berhasil mengotomatisasi sebagian proses perancangan ekonomi gim, pendekatan ini masih beroperasi pada tingkat makro tanpa mempertimbangkan variabilitas preferensi dan strategi pemain dalam skala mikro.

Secara keseluruhan, penelitian terdahulu menunjukkan bahwa algoritma evolusioner dan simulasi berbasis agen non-learning merupakan pendekatan yang efektif dalam optimasi sistem reward dan ekonomi gim. Namun, masih terdapat celah penelitian dalam hal integrasi antara adaptivitas perilaku pemain dengan optimasi algoritmik secara evolusioner pada sistem ekonomi gim yang dinamis.

% Keseimbangan sistem \textit{reward} dalam ekonomi gim merupakan faktor penting dalam menjaga kualitas pengalaman bermain. Sistem \textit{reward} yang tidak seimbang dapat menyebabkan pemain kehilangan motivasi, mengalami kebosanan, atau merasa frustrasi, sehingga berdampak pada keberlanjutan permainan \citep{Holmberg2021}. Desain ekonomi gim, termasuk mekanisme pemberian hadiah, progresi level, dan distribusi sumber daya virtual, berpengaruh langsung terhadap dinamika permainan dan kepuasan pemain.

% Kompleksitas ekonomi gim modern menyebabkan proses penyeimbangan menjadi sulit dilakukan secara manual. Perubahan kecil pada parameter ekonomi sering kali menimbulkan dampak signifikan dan tidak terduga terhadap keseluruhan sistem. Oleh karena itu, berbagai penelitian mengusulkan penggunaan algoritma evolusioner untuk mengotomatisasi proses optimasi dan balancing sistem \textit{reward} \citep{Volz2016}.

% \citet{Volz2016} mendemonstrasikan kelayakan \textit{automatic game balancing} melalui optimasi parameter kartu pada permainan kartu menggunakan algoritma evolusioner berbasis simulasi. Pendekatan ini mampu menghasilkan konfigurasi parameter yang seimbang tanpa memerlukan penalaan manual yang ekstensif. Selanjutnya, \citet{deMesentierSilva2019} menerapkan algoritma evolusioner untuk menyeimbangkan meta permainan \textit{Hearthstone} dengan mengevolusi atribut kartu, sehingga tidak ada strategi yang terlalu dominan.

% Pada \textit{genre} lain, \citeauthor{Gravina2016} memperkenalkan \textit{Constraint Surprise Search} untuk menghasilkan senjata yang beragam namun tetap seimbang dalam gim first person shooter. Sementara itu, Morosan dan Poli \citep{Morosan2017} menunjukkan bahwa algoritma evolusioner dapat digunakan untuk menyeimbangkan unit dalam gim \textit{real time strategy} dengan menyesuaikan parameter unit secara otomatis.

% Sebagian besar penelitian awal berfokus pada komponen permainan secara individual. Namun, \citeauthor{Rogers2023} memperluas pendekatan ini dengan mengoptimasi ekonomi gim secara keseluruhan. Mereka merepresentasikan ekonomi gim sebagai graf alur sumber daya dan menggunakan algoritma evolusioner untuk menargetkan tingkat kompleksitas ekonomi tertentu. Pendekatan ini memungkinkan kontrol yang lebih terstruktur terhadap sistem \textit{reward} dan ekonomi gim.

% Kerangka kerja \textit{Game Economy Evolution} (GEEvo) selanjutnya membagi proses perancangan ekonomi gim menjadi tahap generasi struktur dan tahap optimasi parameter. Algoritma evolusioner digunakan untuk menyesuaikan parameter ekonomi berdasarkan tujuan \textit{balancing} yang ditetapkan, sehingga proses iterasi desain dapat dilakukan secara sistematis dan fleksibel \citep{Rupp_2024}.

% Selain itu, simulasi berbasis agen non \textit{learning} digunakan untuk mengevaluasi sistem \textit{reward} terhadap perilaku pemain yang beragam. Agen direpresentasikan melalui aturan dan heuristik tertentu yang mencerminkan berbagai gaya bermain. Dengan menjalankan simulasi berulang, dampak sistem \textit{reward} terhadap progresi pemain dan stabilitas ekonomi dapat dianalisis secara komprehensif \citep{LaraCabrera2014}.

% Secara keseluruhan, penelitian terdahulu menunjukkan bahwa algoritma evolusioner dan simulasi berbasis agen non \textit{learning} merupakan pendekatan yang efektif dalam optimasi sistem \textit{reward} dan ekonomi gim. 
% Pendekatan ini menyediakan kerangka otomatis dan sistematis untuk mencapai keseimbangan permainan tanpa ketergantungan pada metode pembelajaran mesin. Tabel berikut menyajikan ringkasan dari penjelasan dengan tujuan mengidentifikasi celah penelitian yang masih perlu ditelusuri.

% Definisi simbol centang dan silang
% \newcommand{\cmark}{\ding{51}} % centang
% \newcommand{\xmark}{\ding{55}} % silang

% \begin{table}[htbp]
% \centering
% \small
% \caption{Matriks Research Gap Penelitian Terdahulu}
% \label{tab:research_gap}
% \resizebox{\textwidth}{!}{
% \begin{tabular}{|c|l|c|l|c|c|c|}
% \hline
% \textbf{No} & \textbf{Peneliti} & \textbf{Tahun} & \textbf{Fokus Penelitian} & 
% \textbf{\begin{tabular}[c]{@{}c@{}}Integrasi\\Makro-Mikro\end{tabular}} & 
% \textbf{\begin{tabular}[c]{@{}c@{}}Variasi\\Perilaku\\Pemain\end{tabular}} & 
% \textbf{\begin{tabular}[c]{@{}c@{}}Sistem\\Reward\\Komprehensif\end{tabular}} \\
% \hline
% 1 & Volz et al. & 2016 & 
% \begin{tabular}[c]{@{}l@{}}Automatic game balancing\\pada permainan kartu\end{tabular} & 
% \xmark & \xmark & \xmark \\
% \hline
% 2 & Gravina et al. & 2016 & 
% \begin{tabular}[c]{@{}l@{}}Generasi senjata beragam\\dan seimbang dalam FPS\end{tabular} & 
% \xmark & \xmark & \xmark \\
% \hline
% 3 & Morosan \& Poli & 2017 & 
% \begin{tabular}[c]{@{}l@{}}Balancing unit dalam\\RTS game\end{tabular} & 
% \xmark & \xmark & \xmark \\
% \hline
% 4 & 
% \begin{tabular}[c]{@{}l@{}}de Mesentier\\Silva et al.\end{tabular} & 2019 & 
% \begin{tabular}[c]{@{}l@{}}Balancing meta game\\Hearthstone\end{tabular} & 
% \xmark & \xmark & \xmark \\
% \hline
% 5 & 
% \begin{tabular}[c]{@{}l@{}}Lara-Cabrera\\et al.\end{tabular} & 2014 & 
% \begin{tabular}[c]{@{}l@{}}Evaluasi sistem reward\\menggunakan simulasi agen\end{tabular} & 
% \xmark & \xmark & \cmark \\
% \hline
% 6 & 
% \begin{tabular}[c]{@{}l@{}}Garc\'ia-S\'anchez\\et al.\end{tabular} & 2020 & 
% \begin{tabular}[c]{@{}l@{}}Optimasi agen Hearthstone\\menggunakan algoritma\\evolusioner\end{tabular} & 
% \xmark & \xmark & \xmark \\
% \hline
% 7 & Rogers et al. & 2023 & 
% \begin{tabular}[c]{@{}l@{}}Optimasi ekonomi gim\\menggunakan representasi\\graf\end{tabular} & 
% \cmark & \xmark & \cmark \\
% \hline
% 8 & Rupp \& Eckert & 2024 & 
% \begin{tabular}[c]{@{}l@{}}Generasi dan balancing\\ekonomi gim (GEEvo)\end{tabular} & 
% \cmark & \xmark & \cmark \\
% \hline
% \end{tabular}
% }
% \end{table}


\section{Gim}\label{}
Istilah gim merupakan padanan bahasa Indonesia dari kata \textit{game} yang berarti “permainan” dalam konteks digital maupun non-digital \citep{danny2021valorant}. Dalam KBBI, gim juga didefinisikan sebagai “permainan", terutama yang menggunakan komputer, konsol, atau gawai \citep{KBBI2026}. Hal ini menandakan bahwa istilah gim telah diakui secara linguistik dan sesuai dengan upaya pembakuan bahasa Indonesia terhadap terminologi teknologi dan digital modern.

Dalam kajian ilmiah, gim atau game didefinisikan sebagai sistem interaktif yang melibatkan pemain untuk mencapai tujuan tertentu di bawah seperangkat aturan \citep{salen2004rules}. Gim tidak hanya berfungsi sebagai media hiburan, tetapi juga sebagai sistem kompleks yang mencerminkan dinamika ekonomi, perilaku pengguna, serta sistem penghargaan atau \textit{reward system} yang memengaruhi keterlibatan pemain \citep{hamari2017service}. 

\section{Desain Gim}\label{}

Desain gim merupakan proses sistematis untuk menciptakan pengalaman interaktif yang bermakna melalui kombinasi antara mekanika permainan (\textit{game mechanics}), dinamika, dan estetika. Menurut \cite{schreiber2021game}, desain gim mencakup seluruh keputusan konseptual yang membentuk bagaimana pemain berinteraksi dengan sistem permainan dan bagaimana sistem tersebut memberikan umpan balik. Proses desain ini melibatkan pemahaman mendalam tentang motivasi pemain, keseimbangan sistem ekonomi, serta penyusunan aturan yang mengarahkan perilaku dalam dunia gim.

\section{Ekonomi Gim}\label{}

Ekonomi gim merupakan salah satu aspek dalam desain gim yang mengatur aliran sumber daya virtual di dalam suatu ekosistem gim, seperti mata uang, pengalaman (XP), serta energi \citep{Lin2025}. Pada dasarnya, sistem ekonomi gim bertujuan untuk menyeimbangkan usaha pemain (\textit{player action}) dengan hasil yang diperoleh, sehingga tercipta suatu gim yang adil, menarik, dan berkelanjutan \citep{BowmanKeeneNajera2021}.

Menurut \citet{RanandehMirzaBabaei2023}, komponen utama dalam desain ekonomi gim yaitu :

\begin{enumerate}
    \item Sumber atau \textit{source}, mekanisme yang menambahkan sumber daya ke dalam sistem pemain, seperti hadiah dari misi, \textit{loot box}, atau pembelian dengan uang nyata.
    
    \item Pengeluaran atau \textit{sinks}, mekanisme yang mengeluarkan atau menghapus sumber daya dari sistem pemain, seperti biaya \textit{upgrade}, biaya perbaikan, ataupun konsumsi energi.
    
    \item Konverter atau \textit{converter}, mekanisme yang mengubah suatu jenis sumber daya menjadi jenis lain, seperti \textit{crafting} atau \textit{trading} antar item.
    
    \item Penyimpanan atau \textit{inventory}, Sistem penyimpanan yang mengatur berapa banyak sumber daya yang bisa disimpan atau dimiliki pemain, di dalam dunia gim biasa dibuat dalam bentuk tas, lemari, atau peti.
    
    \item Pedagang atau \textit{Traders}, entitas yang memfasilitasi pertukaran sumber daya pemain dengan pemain atau pemain dengan sistem, seperti NPC shop atau pasar pemain.
    
\end{enumerate}

\begin{figure}[htbp]
\centering
\resizebox{0.7\textwidth}{!}{
\begin{tikzpicture}[
    node distance=3cm and 4cm,
    % Style untuk Source nodes (hijau gelap dengan teks putih)
    source/.style={circle, draw=green!40!black, fill=green!30!black, 
                   text=white, minimum size=1.5cm, thick, font=\small\bfseries},
    % Style untuk Pool nodes (biru gelap dengan teks putih)
    pool/.style={circle, draw=blue!50!black, fill=blue!40!black, 
                 text=white, minimum size=1.6cm, thick, font=\small\bfseries},
    % Style untuk Converter nodes (oranye gelap dengan teks putih)
    converter/.style={circle, draw=orange!50!black, fill=orange!50!black, 
                      text=white, minimum size=1.5cm, thick, font=\small\bfseries},
    % Style untuk Sink nodes (merah gelap dengan teks putih)
    sink/.style={circle, draw=red!50!black, fill=red!40!black, 
                 text=white, minimum size=1.5cm, thick, font=\small\bfseries},
    % Style untuk edge
    arrow/.style={-{Stealth[length=3mm]}, thick, gray!70},
    % Style untuk label pada edge
    edgelabel/.style={font=\small\bfseries, fill=white, inner sep=2pt}
]

% SOURCE Nodes (kiri) - spasi vertikal diperbesar
\node[source] (quest) at (0,2) {\begin{tabular}{c}Quest\\Complete\end{tabular}};
\node[source] (daily) at (0,-1) {Daily Login};
\node[source] (monster) at (0,-4) {\begin{tabular}{c}Monster\\Drop\end{tabular}};

% POOL Nodes (tengah-kiri) - posisi disesuaikan
\node[pool] (xp) at (5,3) {\begin{tabular}{c}XP\\Pool\end{tabular}};
\node[pool] (gold) at (5,0) {\begin{tabular}{c}Gold\\Pool\end{tabular}};
\node[pool] (item) at (5,-3.5) {\begin{tabular}{c}Item\\Pool\end{tabular}};

% SINK Nodes (tengah-kanan atas) - spasi diperbesar
\node[sink] (levelup) at (10,3.5) {\begin{tabular}{c}Level\\Up\end{tabular}};
\node[sink] (upgrade) at (10,1.5) {\begin{tabular}{c}Upgrade\\System\end{tabular}};
\node[sink] (shop) at (10,-0.5) {Shop};

% CONVERTER Node (tengah-kanan bawah) - digeser ke bawah
\node[converter] (crafting) at (10,-5) {Crafting};

% POOL Node (kanan) - digeser
\node[pool] (equipment) at (14,-5) {\begin{tabular}{c}Equipment\\Pool\end{tabular}};

% SINK Node (kanan atas) - digeser
\node[sink] (skill) at (14,2.5) {\begin{tabular}{c}Skill\\Unlock\end{tabular}};

% Edges dengan pembobotan
% Dari Quest Complete
\draw[arrow] (quest) to[out=20,in=160] node[edgelabel, above] {+100} (xp);
\draw[arrow] (quest) to[out=-10,in=160] node[edgelabel, above] {+50} (gold);

% Dari Daily Login
\draw[arrow] (daily) to[out=0,in=180] node[edgelabel, above] {+20} (gold);

% Dari Monster Drop
\draw[arrow] (monster) to[out=60,in=200] node[edgelabel, right] {+10 to +30} (gold);
\draw[arrow] (monster) to[out=0,in=180] node[edgelabel, below] {0.3} (item);

% Dari Gold Pool
\draw[arrow] (gold) to[out=40,in=180] node[edgelabel, above] {-50} (upgrade);
\draw[arrow] (gold) to[out=-20,in=180] node[edgelabel, above] {-100} (shop);

% Dari XP Pool
\draw[arrow] (xp) to[out=0,in=160] node[edgelabel, above] {-500} (levelup);

% Dari Shop ke Item Pool
\draw[arrow] (shop) to[out=220,in=40] node[edgelabel, right] {+1} (item);

% Dari Item Pool ke Crafting
\draw[arrow] (item) to[out=-20,in=180] node[edgelabel, above] {-3} (crafting);

% Dari Crafting ke Equipment Pool
\draw[arrow] (crafting) -- node[edgelabel, above] {+1} (equipment);

% Dari Upgrade System ke Skill Unlock
\draw[arrow] (upgrade) to[out=30,in=180] node[edgelabel, above] {-1} (skill);

% Dari Level Up ke Skill Unlock
\draw[arrow] (levelup) to[out=0,in=140] node[edgelabel, above] {Auto} (skill);

\end{tikzpicture}
}
% Legend
\centering
\begin{tikzpicture}[
    legend/.style={minimum width=0.6cm, minimum height=0.6cm, draw, thick}
]
    \node[legend, fill=green!30!black, draw=green!40!black] (l1) at (0,0) {};
    \node[right=0.2cm of l1, font=\small] {Source (Sumber Reward)};
    
    \node[legend, fill=blue!40!black, draw=blue!50!black] (l2) at (6,0) {};
    \node[right=0.2cm of l2, font=\small] {Pool (Penyimpanan)};
    
    \node[legend, fill=orange!50!black, draw=orange!50!black] (l3) at (0,-0.7) {};
    \node[right=0.2cm of l3, font=\small] {Converter (Konversi)};
    
    \node[legend, fill=red!40!black, draw=red!50!black] (l4) at (6,-0.7) {};
    \node[right=0.2cm of l4, font=\small] {Sink (Pengeluaran)};
\end{tikzpicture}
\caption*{\textbf{Keterangan:} Angka pada panah menunjukkan bobot aliran sumber daya. 
Tanda + berarti penambahan, - berarti pengurangan, desimal (0.3) menunjukkan probabilitas 30\%.}
\caption{Diagram Aliran Ekonomi Gim (diadaptasi dari \cite{Rupp_2024})}
\label{fig:game_economy}
\end{figure}

Keseimbangan antar komponen utama inilah yang menjadi kunci utama keseimbangan sisem ekonomi gim. Jika salah satu komponen utama tidak seimbang (terlalu banyak atau terlalu sedikit), maka permainan akan menjadi terasa sulit atau justru bisa membosankan, sehingga pemain akan kehilangan motivasi untuk bermain gim. Oleh karena itu, desainer ekonomi gim perlu membuat mekanisme yang mampu merespon perubahan perilaku pemain \citep{BowmanKeeneNajera2021}.

\section{Sistem \textit{Reward}}\label{}
Sistem \textit{reward} dapat dimaknai sebagai motivasi pemain atau sebagai kompromi untuk mengurangi kekecewaan \citep{Hamari2010}. Oleh karena itu, dalam desain ekonomi gim, sistem \textit{reward} ditempatkan sebagai komponen teknis yang berfungsi sebagai sumber (\textit{source}) . Sistem ini bertanggung jawab dalam memasukkan atau menambahkan sumber daya virtual ke dalam ekosistem pemain \citep{RanandehMirzaBabaei2023}. Implementasinya dalam berbagai cara mulai dari \textit{item} material atas hasil penyelesaian misi, poin pengalaman kenaikan level, hingga mekanisme crafting yang mengubah input menjadi output \citep{Hamari2010}. 

Desain sistem reward memiliki tujuan utama untuk mencapai \textit{reward balancing}. \textit{Reward balancing} didefinisikan upaya kuantitatif untuk menyamakan kekuatan atau nilai antar elemen permainan melalui optimasi numerik. Keseimbangan dicapai dengan meminimalkan standar deviasi metrik penilaian seperti kerusakan, kesehatan, atau kemajuan ekonomi \citep{9750161}. Keseimbangan ini secara langsung memengaruhi ekonomi makro dalam gim. Sistem \textit{reward} yang tidak stabil menyebabkan percepatan progres permainan yang tidak wajar dan mengganggu keberlanjutan serta stabilitas ekonomi gim secara makro.

Dalam literatur desain ekonomi dan \textit{balancing} gim, sistem \textit{reward} umumnya 
dimodelkan melalui seperangkat parameter numerik yang mengendalikan intensitas, 
frekuensi, dan ketidakpastian aliran sumber daya ke pemain \citep{Rupp_2024}. 
Parameter-parameter tersebut mencakup besaran \textit{reward}, 
yaitu jumlah sumber daya yang diterima per \textit{event}, frekuensi atau interval 
pemberian \textit{reward} yang menentukan laju akumulasi sumber daya, serta parameter probabilistik 
seperti peluang \textit{drop} item langka atau distribusi hasil \textit{loot box} yang 
mengatur ekspektasi matematis dan variansi \textit{reward} antar pemain \citep{Clark2023}.

\section{\textit{Balancing}}\label{}
Balancing dalam game merujuk pada proses penyesuaian mekanik, aturan, serta atribut elemen permainan agar seluruh pilihan yang tersedia bagi pemain berada dalam kondisi yang adil dan layak digunakan. \cite{RollingsAdams2003} menyatakan bahwa sebuah game dikatakan seimbang apabila keberhasilan pemain ditentukan oleh keterampilan bermain, bukan oleh keberuntungan semata atau dominasi satu strategi tertentu. Keseimbangan ini bertujuan mencegah munculnya \textit{dominant strategy} yang dapat menghilangkan variasi dan makna dalam pengambilan keputusan.

\cite{Sirlin2016} menekankan bahwa balancing tidak berarti seluruh opsi harus memiliki kekuatan yang sama, melainkan setiap opsi harus memiliki kelebihan dan kelemahan yang jelas sehingga menghasilkan pilihan yang bermakna dan adil, terutama pada permainan kompetitif. Opsi yang terlalu kuat akan mendominasi, sedangkan opsi yang terlalu lemah menjadi tidak relevan dan mengurangi kedalaman permainan.

Novak serta Rollings dan Adams membedakan balancing menjadi dua pendekatan, yaitu \textit{static balancing} dan \textit{dynamic balancing}. Static balancing dilakukan pada tahap perancangan dengan mengatur nilai dasar seperti kekuatan, biaya, dan rasio risiko-imbalan, sedangkan dynamic balancing terjadi selama permainan berlangsung melalui penyesuaian tingkat kesulitan atau mekanik untuk menjaga tantangan tetap seimbang bagi pemain dengan tingkat kemampuan berbeda \cite{Novak2012}.

% Portnow memperkenalkan konsep \textit{perfect imbalance}, yaitu gagasan bahwa perbedaan kekuatan yang kecil antar elemen justru dapat meningkatkan dinamika strategi dan mendorong terbentuknya metagame. Selama ketimpangan tersebut tidak menciptakan elemen yang benar-benar mendominasi, variasi strategi tetap terjaga dan pengalaman bermain menjadi lebih menarik \cite{Portnow2013}. Dengan demikian, balancing dipahami bukan sekadar penyamaan kekuatan, tetapi sebagai upaya menjaga keadilan, kedalaman strategi, serta kebermaknaan setiap pilihan dalam permainan.


\section{Graf Berarah}\label{}

Dalam konteks ekonomi gim, graf berarah (\textit{directed graph}) digunakan untuk merepresentasikan aliran sumber daya antar entitas dalam sistem. Graf ini memungkinkan analisis struktural terhadap hubungan antara sumber (\textit{sources}), pengeluaran (\textit{sinks}), serta konversi sumber daya dalam ekosistem gim. Representasi matematis dari ekonomi gim dalam bentuk graf berarah membantu dalam memodelkan keterkaitan dinamis antara komponen-komponen sistem \textit{reward} dan ekonomi virtual secara kuantitatif \citep{Rogers2023}.

Sebuah graf berarah $G$ didefinisikan sebagai pasangan terurut:
\begin{equation}
    G = (V, E),
\end{equation}
dengan $V = \{v_1, v_2, \dots, v_n\}$ merupakan himpunan \textit{vertex} (node) yang merepresentasikan komponen ekonomi seperti penyimpanan, toko, atau sumber daya, dan $E \subseteq V \times V$ adalah himpunan \textit{edge} berarah $(v_i, v_j)$ yang menyatakan adanya aliran sumber daya dari $v_i$ menuju $v_j$.

Bobot pada setiap sisi dilambangkan sebagai $w_{ij}$ yang menggambarkan besar aliran sumber daya dari node $i$ ke node $j$. Secara matematis, hubungan antar node dapat dituliskan dalam bentuk matriks bobot:
\begin{equation}
    W = [w_{ij}]_{n \times n}, \quad w_{ij} \ge 0,
\end{equation}
Jika $w_{ij} = 0$, maka tidak ada aliran langsung dari $v_i$ ke $v_j$. Jumlah aliran keluar (\textit{outflow}) dan masuk (\textit{inflow}) dari suatu node dapat dihitung sebagai:
\begin{align}
    \text{Outflow}(v_i) &= \sum_{j=1}^{n} w_{ij}, \\
    \text{Inflow}(v_i) &= \sum_{j=1}^{n} w_{ji}.
\end{align}

\citet{Rogers2023} menunjukkan bahwa representasi ekonomi gim dalam bentuk graf berarah memungkinkan penerapan algoritma evolusioner untuk mengoptimalkan kompleksitas dan kestabilan sistem ekonomi secara adaptif. Struktur graf dapat dioptimasi dengan fungsi objektif yang meminimalkan variansi ketidakseimbangan aliran:
\begin{equation}
    f(G) = \sum_{i=1}^{n} \left( \text{Inflow}(v_i) - \text{Outflow}(v_i) \right)^2.
\end{equation}
Tujuan algoritma evolusioner dalam konteks ini adalah menemukan konfigurasi graf yang meminimalkan $f(G)$, sehingga sistem \textit{reward} menjadi stabil dan seimbang.

\citet{Rupp_2024} memperkenalkan kerangka \textit{GEEvo} (\textit{Game Economy Evolution}) yang menggunakan representasi graf berarah untuk menghasilkan dan menyeimbangkan ekonomi gim secara otomatis. Node dan edge direpresentasikan sebagai gen dalam kromosom evolusioner, sementara operasi mutasi dan \textit{crossover} digunakan untuk mengeksplorasi struktur graf yang lebih optimal terhadap tujuan keseimbangan ekonomi gim. \citet{Lin2025} juga menegaskan bahwa penggunaan model graf memungkinkan pengembang untuk mensimulasikan dampak perubahan nilai virtual \textit{currency} terhadap perilaku pemain, sehingga sistem ekonomi dapat dirancang secara lebih adaptif terhadap dinamika pengguna.

Secara keseluruhan, pendekatan berbasis graf berarah memberikan dasar matematis dan komputasional yang kuat untuk memahami serta mengoptimasi aliran sumber daya dalam ekonomi gim. Penggabungan model ini dengan algoritma evolusioner memungkinkan terwujudnya sistem \textit{reward} yang tidak hanya seimbang, tetapi juga adaptif terhadap dinamika permainan dan perilaku pemain  \citep{Rogers2023}.


\section{Motivasi Pemain}
Tingkat keterlibatan, retensi, serta perilaku pemain ditentukan oleh motivasi pemain itu sendiri \citep{Bruhlmann2020}. Dalam desain ekonomi gim, memahami motivasi pemain sangat penting karena menentukan sistem \textit{reward} yang sesuai untuk mendorong perilaku yang diharapkan serta mempertahankan keseimbangan antara tantangan dan kepuasan. 

Secara umum, motivasi pemain dapat dibedakan menjadi dua kelas utama :
\begin{enumerate}
    \item Motivasi intrinsik, dorongan yang muncul dari kepuasan diri pemain, seperti rasa pencapaian, eksplorasi, atau kreativitas \citep{HoegVanDerKaapDeeder2023}.
    \item Motivasi ekstrinsik, dorongan yang muncul dari faktor eksternal, seperti hadiah, pengakuan sosial, atau pencapaian tertentu dalam gim \citep{HoegVanDerKaapDeeder2023}.
\end{enumerate}

dalam \textit{Self Determination Theory} dijelaskan bahwa motivasi manusia didorong oleh tiga kebutuhan psikologis dasar \citep{UysalYildirim2016} :

\begin{itemize}
    \item Autonomi, perasaan bahwa pemain memiliki kendali terhadap tindakan dan pilihannya di dalam gim
    \item Kompetensi, perasaan mampu dan efektif dalam menghadapi tantangan
    \item Keterhubungan, rasa memiliki hubungan sosial yang positif dengan pemain lain atau karakter di dalam gim
\end{itemize}

Selain itu, \textit{Flow Theory} yang diperkenalkan oleh Csikszentmihalyi mengatakan bahwa pemain akan berada pada keadaan \textit{flow} ketika tantangan dan kemampuan mereka berada pada tingkat yang seimbang \citep{2017Inverted}. Dalam kondisi \textit{flow}, pemain akan merasa terlibat penuh dan akan melupakan persepsi terhadap waktu yang merupakan indikator kuat dari keterlibatan optimal. Desain \textit{reward} system harus mampu mengakomodasi ketiga kebutuhan SDT sekaligus menjaga kondisi \textit{flow} agar pemain tetap terlibat \citep{BuilCatalanMartinez2019}.

\section{Algoritma Evolusioner}

Algoritma evolusioner merupakan salah satu pendekatan optimasi berbasis populasi yang terinspirasi dari evolusi biologis seperti seleksi alam, mutasi, dan reproduksi \citep{Herel2022}. Dalam konteks sistem ekonomi gim, algoritma evolusioner dapat digunakan untuk mencari konfigurasi optimal dari parameter sistem \textit{reward} agar tercipta keseimbangan ekonomi gim dan keterlibatan pemain tetap tercapai \citep{MotepalliJacobsen2021}. 

Prinsip dasar algoritma evolusioner dimulai dengan populasi solusi kandidat yang merepresentasikan berbagai kemungkinan sistem \textit{reward}. Setiap individu dalam populasi diatur menggunakan fungsi \textit{fitness} yang menentukan seberapa baik solusi tersebut mencapai tujuan tertentu, seperti stabilitas ekonomi gim, retensi pemain, atau tingkat kepuasan \citep{Bai2024An}. 

\begin{algorithm}[H]
\caption{Kerangka Umum Algoritma Evolusioner untuk Optimasi Sistem \textit{Reward Gim}}
\label{alg:ea-reward}
\begin{algorithmic}[1]
\INPUT  \parbox[t]{\dimexpr\linewidth-\labelwidth}{Populasi awal $P_0$, fungsi \textit{fitness} $f(\cdot)$, probabilitas \textit{crossover} $p_c$, probabilitas mutasi $p_m$, kondisi terminasi}
\OUTPUT Solusi terbaik $x^*$
\State Inisialisasi populasi awal $P_0$
\State Evaluasi nilai \textit{fitness} setiap individu dalam $P_0$
\While{kondisi terminasi belum terpenuhi}
    \State Seleksi individu dari populasi berdasarkan nilai \textit{fitness}
    \State Lakukan operasi \textit{crossover} dengan probabilitas $p_c$
    \State Lakukan operasi mutasi dengan probabilitas $p_m$
    \State Bentuk populasi baru $P_{t+1}$
    \State Evaluasi nilai \textit{fitness} setiap individu dalam $P_{t+1}$
\EndWhile
\State Tentukan individu dengan nilai \textit{fitness} terbaik sebagai $x^*$
\State \Return $x^*$
\end{algorithmic}
\end{algorithm}

Proses pada Algoritma~\ref{alg:ea-reward} diawali dengan inisialisasi populasi awal \(P_0={x_1,x_2,…,x_N}\) yang merepresentasikan berbagai kemungkinan konfigurasi sistem \textit{reward}. Inisialisasi ini dapat dilakukan secara acak maupun berdasarkan parameter ekonomi gim yang telah ada sebelumnya. Setiap individu dalam populasi kemudian dievaluasi menggunakan fungsi \textit{fitness} $f(\cdot)$ yang mengukur kualitas solusi terhadap kriteria tertentu, seperti stabilitas ekonomi gim, keseimbangan antara \textit{reward} dan \textit{sink}, serta tingkat keterlibatan pemain. Perlu ditekankan bahwa fungsi \textit{fitness} dalam algoritma evolusioner tidak bersifat universal dan dapat dirancang secara berbeda-beda, bergantung pada tujuan optimasi dan karakteristik sistem yang dianalisis \citep{Baresel2002Fitness}. Sebagai contoh, pada kerangka \textit{GEEvo (Generate \& Balancing Game Economies)},
fungsi fitness dirancang berbeda pada dua tahap optimasi. Pada tahap \textit{generator},
fungsi fitness dirumuskan sebagai:
\begin{equation}
f_{\text{gen}}(G) = \sum_{v \in V} \mathrm{validate}(v),
\end{equation}
dengan keterangan:
\begin{itemize}
    \item $G$ : graf ekonomi yang dihasilkan,
    \item $V$ : himpunan node pada graf,
    \item $v$ : node individu pada graf,
    \item $\mathrm{validate}(v)$ : Fungsi bantu yang menghitung jumlah batasan yang tidak terpenuhi oleh node $v$.
\end{itemize}

Setelah struktur graf dinyatakan valid, tahap \textit{balancer} menggunakan fungsi
fitness berbasis hasil simulasi yang dirumuskan sebagai:
\begin{equation}
    f_1(s^{p_j}_t, x) = \alpha + \frac{1}{m} \sum_{i=1}^{m} prop(s^{p_j}_{t_i}, x),
\end{equation}

\begin{equation}
    prop(s^{p_j}_t, x) = 
    \begin{cases} 
        s^{p_j}_t \cdot \frac{1}{x}, & \text{if } x > s^{p_j}_t \\
        x \cdot \frac{1}{s^{p_j}_t}, & \text{if } x \le s^{p_j}_t ,
    \end{cases}
\end{equation}
dengan keterangan:
\begin{itemize}
    \item $s^{p_j}_t$: Jumlah sumber daya dalam pool $p_j$ pada langkah waktu $t$.
    \item $x$: Nilai target absolut yang ingin dicapai (misalnya: 100 gold).
    \item $\alpha$: Parameter ambang batas (\textit{threshold}) untuk toleransi rata-rata, membantu algoritma berhenti jika nilai sudah cukup dekat (maksimum fitness adalah $1+\alpha$).
    \item $m$: Jumlah simulasi yang dijalankan untuk memitigasi keacakan (\textit{randomness}).
    \item $i$: Indeks untuk setiap putaran simulasi.
    \item $prop$: Fungsi proporsional untuk menghitung seberapa dekat nilai saat ini dengan target.
\end{itemize}


Tahap seleksi bertujuan untuk memilih individu dengan nilai \textit{fitness} yang lebih baik agar memiliki peluang lebih besar untuk berkontribusi pada generasi berikutnya. Seleksi ini mencerminkan prinsip seleksi alam, di mana solusi yang lebih adaptif terhadap tujuan sistem akan lebih mungkin dipertahankan. Selanjutnya, operasi \textit{crossover} dilakukan dengan mengombinasikan dua individu terpilih untuk menghasilkan solusi baru yang berpotensi mewarisi karakteristik unggul dari keduanya. Untuk menjaga keberagaman populasi dan mencegah konvergensi prematur, operasi mutasi diterapkan dengan probabilitas tertentu. Mutasi dimodelkan dengan menghasilkan perubahan kecil pada parameter solusi, yaitu dengan menambahkan gangguan acak berdistribusi normal pada setiap parameter : 
\begin{equation}
x_i' = x_i + \varepsilon, \qquad \varepsilon \sim \mathcal{N}(0, \sigma^2).
\end{equation}

Populasi baru $P_{t+1}$ kemudian dievaluasi kembali menggunakan fungsi \textit{fitness}, dan proses ini diulang secara iteratif hingga kondisi terminasi terpenuhi, misalnya ketika jumlah generasi maksimum tercapai atau peningkatan nilai \textit{fitness} tidak lagi signifikan. Individu dengan nilai \textit{fitness} terbaik pada akhir proses evolusi dipilih sebagai solusi optimal $x^*$. Proses menyeimbangkan sistem \textit{reward} atau ekonomi gim dengan menggunakan algoritma evolusioner dapat divisualisasikan pada gambar berikut.

\begin{figure}[htbp]
\centering
\begin{tikzpicture}[
    node distance=2cm,
    block/.style={rectangle, draw, thick, minimum width=2.8cm, minimum height=1cm, align=center},
    arrow/.style={->, thick},
    starnode/.style={draw, thick, circle, minimum size=0.6cm}
]

% Draw dashed box
\draw[dashed, thick] (-8, -3.2) rectangle (6, 2);
\node[anchor=north west] at (-6, 2);

% Nodes
\node[block] (fitness) at (0, 0.5) {Fitness $f(s_t, x)$};
\node[block] (update) at (-3, -1.5) {Pembaruan\\Bobot};
\node[block] (economy) at (3, -1.5) {Simulasi\\Ekonomi Gim};

% Star nodes
\node[starnode] (star1) at (-5, -1.5) {$*$};
\node[starnode] (star2) at (0, -1.5) {$*$};

% Labels below star1
\node[align=center, font=\footnotesize] at (-6.7, -1.5) {Bobot Graf\\Ekonomi Gim};

% Arrows
\draw[arrow] (fitness.west) -| node[above, pos=0.25] {$f_t$} (update.north);
\draw[arrow] (economy.north) |- node[above, pos=0.75] {$s_t$} (fitness.east);
\draw[arrow] (star1) -- (update.west);
\draw[arrow] (update.east) -- (star2.west);
\draw[arrow] (star2.east) -- (economy.west);

\end{tikzpicture}
\caption{Struktur Penyeimbangan \citep{Rupp_2024}}
\label{fig:balancer_structure}
\end{figure}

Dengan memanfaatkan algoritma evolusioner, beberapa aspek dalam ekonomi gim dapat dioptimasi seperti frekuensi distribusi \textit{reward}, rasio \textit{reward-sink}, parameter probabilistik \textit{loot table}, serta kurva progresi ekonomi pemain

\section{Simulasi Agen}

Simulasi berbasis agen merupakan simulasi pemodelan yang menggambarkan sistem kompleks melalui interaksi antara entitas mandiri yang disebut agen \citep{MacalNorth2005}. Setiap agen memiliki karakteristik, perilaku, serta kemampuan berinteraksi dengan agen lain maupun dengan lingkungan tempat mereka berada \citep{DatserisVahdatiDuBois2021}. Dalam konteks ekonomi gim, simulasi agen digunakan untuk memodelkan perilaku pemain dan menganalisis bagaimana perubahan dalam sistem \textit{reward} memengaruhi ekonomi gim secara makro. Dengan pendekatan ini, pengembang dapat mengetahui kebijakan dari \textit{reward balancing} tanpa perlu testing secara langsung.

Model simulasi ini dapat digunakan untuk mengamati fenomena seperti \citep{NichollsAmelungStudent2017} :

\begin{itemize}
    \item Distribusi kekayaan virtual, hasil dari kombinasi reward, sink, dan interaksi antar pemain.
    \item Efek keseimbangan reward terhadap retensi agen, yaitu berapa lama agen tetap aktif dalam sistem.
    \item Emergent behavior, seperti inflasi ekonomi gim akibat reward berlebih atau ketimpangan antara pemain aktif dan pasif.
\end{itemize}

Dalam penelitian sistem ekonomi gim, hal ini memungkinkan analisis \textit{equilibrium} dinamis, di mana parameter sistem \textit{reward} dapat disesuaikan secara iteratif hingga tercapai stabilitas.

    \chapter{Metode Penelitian}\label{ch:method}
\section{Tahapan Penelitian}\label{sec:stages}
Bab ini membahas tahapan penelitian dalam pelaksanaan Tugas Akhir. Penelitian ini berfokus pada perancangan model dan optimasi sistem \textit{reward} dalam ekonomi gim menggunakan simulasi berbasis agen dan algoritma evolusioner. Pendekatan yang digunakan bertujuan untuk merepresentasikan aliran sumber daya dalam sistem \textit{reward} serta menganalisis pengaruhnya terhadap dinamika ekonomi gim. Alur tahapan penelitian ini ditunjukkan pada Gambar 3.1.

\setlength{\fboxsep}{10pt} 
\begin{figure}[H]
\centering
\fbox{
\begin{tikzpicture}[
    node distance=0.9cm,
    every node/.style={
        draw,
        rectangle,
        rounded corners,
        align=center,
        minimum width=5cm,
        minimum height=1cm,
        font=\small
    },
    arrow/.style={
        ->,
        thick
    }
]

\node (literature) {Studi Literatur};
\node (analysis) [below=of literature] {Analisis dan Perancangan Model Sistem Reward};
\node (simulation) [below=of analysis] {Analisis dan Perancangan Simulasi Agen};
\node (evolutionary) [below=of simulation] {Perancangan Algoritma Evolusioner};
\node (evaluation) [below=of evolutionary] {Uji Coba dan Diskusi};
\node (conclusion) [below=of evaluation] {Analisis Hasil dan Penarikan Kesimpulan};

\draw[arrow] (literature) -- (analysis);
\draw[arrow] (analysis) -- (simulation);
\draw[arrow] (simulation) -- (evolutionary);
\draw[arrow] (evolutionary) -- (evaluation);
\draw[arrow] (evaluation) -- (conclusion);

\end{tikzpicture}
}
\caption{Tahapan Penelitian.}
\label{fig:tahapan_penelitian}
\end{figure}


Tahapan penelitian ini diawali dengan studi literatur yang membahas konsep ekonomi gim, sistem \textit{reward}, simulasi berbasis agen, dan algoritma evolusioner. Studi literatur ini digunakan sebagai dasar dalam menyusun spesifikasi model sistem \textit{reward} serta metode algoritma yang akan digunakan dalam penelitian. Selanjutnya dilakukan analisis dan spesifikasi model ekonomi gim, termasuk penentuan komponen sistem, hubungan antar komponen, serta representasi model menggunakan graf berarah. Tahap berikutnya adalah perancangan simulasi berbasis agen untuk merepresentasikan perilaku pemain yang berinteraksi di dalam sistem \textit{reward}. Simulasi ini dirancang untuk menghasilkan data yang menggambarkan dinamika sistem \textit{reward} berdasarkan berbagai konfigurasi parameter. Berdasarkan hasil simulasi tersebut, algoritma evolusioner digunakan untuk mengoptimasi parameter sistem \textit{reward} sesuai dengan kriteria evaluasi yang ditentukan. Tahap akhir penelitian adalah uji coba dan diskusi serta evaluasi, yang bertujuan untuk membandingkan kinerja sistem \textit{reward} hasil optimasi dengan konfigurasi awal.

\section{Studi Literatur}
Studi literatur dilakukan guna menunjang penyusunan tugas akhir serta memperdalam pemahaman mengenai \textit{balancing} sistem \textit{reward} ekonomi gim, algoritma evolusioner, dan simulasi agen. Literatur dapat berasal dari jurnal ataupun prosiding konferensi. Jurnal dan prosiding konferensi dapat diperoleh melalui basis data ilmiah seperti \href{https://ieeexplore.ieee.org/Xplore/home.jsp}{IEEE Xplore}, \href{https://link.springer.com/}{SpringerLink}, \href{https://www.sciencedirect.com/}{ScienceDirect}, serta \href{https://arxiv.org/}{arXiv}. Literatur yang dipilih dibatasi pada penelitian yang membahas simulasi dan optimasi sistem dalam konteks gim.


\section{Analisis dan Perancangan Model Sistem \textit{Reward}}
Perancangan model sistem \textit{reward} dilakukan dengan merepresentasikan dinamika aliran sumber daya pada tingkat mikro ke dalam bentuk graf terarah. Setiap node dalam graf memiliki fungsi tersendiri, beberapa diantaranya menghasilkan sumber daya, menyimpan, melakukan distribusi berdasarkan bobot probabilistik, melakukan konversi, serta menghilangkan sumber daya dari ekosistem. Graf dirancang agar memenuhi aturan hubungan antar node sehingga dapat digunakan sebagai dasar simulasi ekonomi gim. Hubungan antar node diwakili oleh edge yang memiliki nilai untuk merepresentasikan parameter \textit{reward} sebagai aliran sumber daya yang akan dioptimasi. Parameter ini mencakup besaran \textit{reward}, frekuensi pemberian \textit{reward}, probabilitas distribusi pada node tertentu, biaya konversi, serta nilai \textit{sink} yang memengaruhi keseimbangan atau stabilitas sistem. Setiap parameter direpresentasikan oleh nilai numerik sehingga dapat dimodifikasi langsung oleh algoritma evolusioner.

\section{Analisis dan Perancangan Simulasi Agen}
Simulasi agen dirancang untuk merepresentasikan pemain dalam berinteraksi di dalam ekosistem gim. Pemain dimodelkan sebagai agen berbasis aturan (\textit{ruled-based}) yang diberi atribut perilaku dan dipengaruhi oleh preferensi serta kecenderungan untuk melakukan suatu aksi. Simulasi dilakukan dalam beberapa langkah waktu, di mana setiap agen melakukan aksi seperti memperoleh \textit{reward}, melakukan \textit{crafting}, atau menghabiskan sumber daya. Perubahan atas hasil aksi agen direkam pada graf ekonomi dan diolah kembali untuk iterasi berikutnya.

Agen dalam simulasi tidak diasumsikan bersifat homogen, 
melainkan memiliki variasi perilaku yang 
merepresentasikan perbedaan gaya bermain pemain. 
Variasi ini memengaruhi probabilitas agen dalam 
memilih aksi ekonomi, seperti menyelesaikan quest, 
melakukan crafting, berbelanja, atau melakukan upgrade.

\section{Perancangan Algoritma Evolusioner}
Perumusan algoritma evolusioner dilakukan untuk mengoptimalisasi parameter \textit{reward} atas hasil dari simulasi agen. Setiap individu dalam populasi mewakili satu konfigurasi parameter. Evaluasi dilakukan dengan mensimulasikan ekonomi gim dengan iterasi tertentu dan kemudian dihitung menggunakan fungsi fitness berdasarkan perbedaan antara kondisi ekonomi aktual dan target yang diharapkan. Proses seleksi, crossover, dan mutasi digunakan untuk menghasilkan populasi baru yang berpotensi menciptakan konfigurasi yang lebih baik. 

Integrasi antara simulasi agen dan algoritma evolusioner dilakukan dengan menghubungkan parameter yang dihasilkan oleh algoritma evolusioner dengan simulasi agen. Setiap individu dievaluasi melalui simulasi agen sehingga diperoleh metrik ekonomi yang diperlukan untuk menghitung \textit{fitness}. Fungsi \textit{fitness} dalam penelitian ini akan diadaptasi dari kerangka GEEvo 
dengan modifikasi untuk mengakomodasi hasil simulasi agen

Evaluasi sistem reward dilakukan secara terintegrasi
melalui fungsi fitness dalam algoritma evolusioner. Tujuan penyeimbangan
ekonomi gim direpresentasikan langsung oleh fungsi fitness, yang
mengukur kualitas suatu konfigurasi parameter sistem reward berdasarkan
hasil simulasi ekonomi gim.

\section{Uji Coba dan Diskusi}
Pada tahap uji coba dan diskusi, eksperimen dilakukan melalui dua tipe uji coba, yaitu uji coba menggunakan model ekonomi gim yang dibangkitkan secara generatif dan uji coba menggunakan model ekonomi gim yang telah ada sebagai data eksternal. Studi kasus untuk masukan eksternal atau model ekonomi gim yang telah ada berjumlah 2.
Data masukan eksperimen yang dibangkitkan berupa konfigurasi struktur ekonomi, parameter sistem \textit{reward}, parameter teknis simulasi, serta tipe agen. Data masukan ini ditentukan sebelum simulasi dijalankan dan berfungsi sebagai kondisi awal eksperimen.

Konfigurasi struktur ekonomi meliputi ukuran graf ekonomi secara keseluruhan serta komposisi jenis node yang digunakan, seperti node sumber (\textit{source}), node probabilistik (\textit{random gate}), node konversi sumber daya (\textit{convert}), node penyimpanan (\textit{pool}), dan node pembuangan (\textit{sinks}), yang bersama-sama menentukan kompleksitas aliran sumber daya dalam simulasi. Selain struktur ekonomi, data masukan juga mencakup parameter sistem \textit{reward} yang menjadi target optimasi, seperti nilai target sumber daya yang ingin dicapai pada \textit{pool} tertentu dalam jangka waktu simulasi tertentu serta parameter toleransi keseimbangan untuk mengatur tingkat presisi keseimbangan ekonomi, sehingga sistem tidak hanya mencapai target absolut tetapi juga mampu mempertahankan stabilitas aliran sumber daya.

Eksperimen dilakukan dengan menjalankan beberapa skenario yang merepresentasikan variasi perilaku pemain pada model ekonomi. Setiap skenario dievaluasi menggunakan algoritma evolusioner untuk memperoleh konfigurasi yang optimal. Hasil data keluaran berupa deret waktu jumlah sumber daya pada setiap pool ekonomi gim serta nilai fungsi fitness pada setiap generasi algoritma evolusioner. Hasil eksperimen kemudian dibandingkan dengan kondisi awal untuk melihat peningkatan \textit{balancing} ekonomi serta dampak perubahan parameter terhadap perilaku agen. Proses ini dilakukan untuk validasi efektivitas model dalam menjaga sistem reward tetap \textit{balance}. Evaluasi dilakukan menggunakan metrik keseimbangan 
aliran sumber daya pada pool serta nilai fitness yang 
dihasilkan oleh algoritma evolusioner.

\section{Analisis Hasil dan Penarikan Kesimpulan}
Setelah eksperimen dilakukan, data simulasi dianalisis untuk menarik kesimpulan keberhasilan model mengoptimasi parameter. Analisis mencakup kestabilan sumber daya, keseimbangan, serta perubahan perilaku sistem sebelum dan sesudah optimasi.


% \section{Pendekatan Penelitian}
% Penelitian ini menggunakan pendekatan komputasional-eksperimental, di mana model ekonomi gim direpresentasikan dalam bentuk graf terarah dan dianalisis menggunakan algoritma evolusioner. Pendekatan ini dipilih karena masalah sistem \textit{reward balancing} yang bersifat non-linier, kompleks, dan melibatkan banyak parameter yang saling berinteraksi, sehingga sulit diselesaikan dengan metode analitik. 
% Simulasi agen dilakukan untuk merepresentasikan eksperimen yang dilakukan oleh pemain dan dinamika aliran sumber daya dalam gim. Hasil simulasi menjadi dasar evaluasi fungsi \textit{fitness} untuk proses evolusi.

% \section{Desain Sistem Penelitian}
% Penelitian ini dibagi menjadi empat tahapan utama :
% \subsection{Perancangan Struktur Ekonomi Gim Berbasis Agen}
% Sistem ekonomi gim direpresentasikan sebagai graf terarah yang memiliki node dan edge. Node mewakili \textit{source, pool, converter, random gate,} dan \textit{drain}. 
% % \textit{Source} berfungsi untuk menciptakan sumber daya dan memasukkannya ke dalam sistem ekonomi, \textit{pool} sama seperti \textit{inventory} untuk menyimpan sumber daya ke dalam ekosistem pemain, \textit{converter} berfungsi untuk mengubah jenis sumber daya, \textit{random gate} berfungsi mendistribusikan sumberdaya berdasarkan bobot probabilitas, \textit{drain} merupakan nama lain dari sinks yang berfungsi menghapus sumber daya dari sistem. 
% \subsection{Formulasi Parameter Reward yang Akan Dioptimasi}
% Beberapa parameter reward yang akan dioptimasi diantaranya :
% \begin{itemize}
%     \item frekuensi \textit{reward},
%     \item distribusi probabilitas loot,
%     \item biaya \textit{upgrade (sink)},
%     \item konversi \textit{crafting},
%     \item progresi ekonomi pemain.
% \end{itemize}

% \subsection{Simulasi Agen untuk Mengukur Dampak \textit{Reward}}
% Setiap iterasi menghasilkan data berupa :
% \begin{itemize}
%     \item distribusi kekayaan agen,
%     \item progresi pemain,
%     \item retensi agen,
%     \item stabilitas ekonomi makro.
% \end{itemize}
% \subsection{Optimasi dengan Algoritma Evolusioner}
% Digunakan algoritma evolusioner untuk mencari konfigurasi terbaik berdasarkan \textit{multi-objective fitness}
% \section{Model Ekonomi Gim}
% Model ekonomi gim direpresentasikan dalam sistem graf terarah, di mana setiap elemen ekonomi gim dimodelkan sebagai node dengan perilaku spesifik.

% \[
% G = (V,E)
% \]
% dengan,
% \begin{itemize}
%     \item \(V\): node yang merepresntasikan elemen ekonomi gim
%     \item \(E\): edge yang merepresentasikan aliran sumber daya
% \end{itemize}
% \subsection{Representasi Node}
% Beberapa definisi node yang digunakan dalam model ekonomi gim ini :
% \begin{itemize}
%     \item \textit{Source} : berfungsi untuk menciptakan sumber daya dan memasukkannya ke dalam sistem ekonomi.
%     \item \textit{Random Gate} : berfungsi mendistribusikan sumberdaya berdasarkan bobot probabilitas.
%     \item \textit{Pool} : sama seperti \textit{inventory} untuk menyimpan sumber daya ke dalam ekosistem pemain.
%     \item \textit{Converter} : berfungsi untuk mengubah jenis sumber daya.
%     \item \textit{Drain} : merupakan nama lain dari sinks yang berfungsi menghapus sumber daya dari sistem.
% \end{itemize}

% \section{Perancangan Agen}
% Setiap agen merepresentasikan pemain yang memiliki atribut :
% \begin{itemize}
%     \item tingkat aktivitas,
%     \item preferensi reward,
%     \item probabilitas interaksi,
%     \item retensi dan durasi bermain. 
% \end{itemize}
% Dalam setiap simulasi, agen perilaku yang akan dilakukan agen diantaranya :
% \begin{itemize}
%     \item melakukan aksi (\textit{quest, crafting, upgrade}),
%     \item menerima reward,
%     \item menghabiskan sumber daya (\textit{sink}),
%     \item memperbarui status progres dan engagement.
% \end{itemize}
% Perilaku yang dilakukan agen sederhana tetapi cukup untuk merepresntasikan dinamika ekonomi gim yang kompleks.

% \section{Algoritma Evolusioner Untuk Optimasi Reward}
% \subsection{Representasi Individu}
% Setiap individu merepresentasikan kombinasi parameter \textit{reward}, seperti :
% \begin{itemize}
%     \item jumlah reward,
%     \item probabilitas random gate,
%     \item biaya crafting,
%     \item rasio \textit{sink-reward},
%     \item interval pemberian \textit{reward}
% \end{itemize}
% Semua parameter disimpan sebagai vektor bilangan real positif.

% \subsection{Inisialisasi Populasi}
% Populasi awal dibuat secara acak dengan ukuran N.
% Setiap individu kemudian diuji melalui simulasi untuk mendapatkan nilai fitness awal.

% \subsection{Evaluasi \textit{Fitness}}
% \textit{Fitness} didefinisikan berdasarkan kemampuan gim mencapai keseimbangan.
% \subsubsection{Fitness Menuju Target Reward}
% Fungsi fitness utama yang digunakan :
% \[
% f_1(s_t^{pj},x)= \alpha+\frac{1}{m}\sum_{i=1}^m prop(s_t^{pj},x)
% \]
% dengan fungsi proporsional:

% \[
% prop(s_t^{p_j}, x) =
% \begin{cases}
% \frac{s_t^{p_j}}{x}, & x>0 \text{ jika } x>s_t^{p_j} \\[6pt]
% \frac{x}{s_t^{p_j}}, & s_t^{p_j}>0 \text{ jika } x \le s_t^{p_j}
% \end{cases}
% \]

% Interpretasi :
% \begin{itemize}
%     \item Nilai fitness terbaik = \(\alpha+1\)
%     \item semakin mendekati \(\alpha+1\), semakin seimbang parameter reward
% \end{itemize}
% \subsubsection{Fitness Untuk Dua Ekonomi Sekaligus}
% Jika dua ekonomi harus diseimbangkan:
% \[
% f_2(s_t^{p_j},s_t^{p_k})=f_1(s_t^{p_j},s_t^{p_k})
% \]
% Digunakan untuk perbandingan kelas atau sistem reward berbeda.

% \subsection{Seleksi}
% Metode seleksi seperti tournament selection atau roulette wheel digunakan untuk memilih individu dengan fitness terbaik

% \subsection{\textit{Crossover}}
% Parameter dari dua individu digabungkan menggunakan mekanisme rekombinasi untuk menghasilkan keturunan baru.

% \subsection{Mutasi}
% Mutasi dilakukan dengan :
% \begin{itemize}
%     \item menambahkan \textit{noise} kecil ke parameter,
%     \item mengubah probabilitas \textit{reward},
%     \item mengubah frekuensi distribusi.
% \end{itemize}

% \subsection{Terminasi}
% Algoritma berhenti jika memenuhi kondisi :
% \begin{itemize}
%     \item jumlah generasi maksimum tercapai, atau
%     \item nilai fitness stabil dan tidak lagi meningkat signifikan.
% \end{itemize}
    \chapter{Hasil dan Pembahasan}\label{ch:result}

Bab ini membahas hasil implementasi dan pengujian dari \textit{balancing} sistem \textit{reward} yang telah dirancang pada bab sebelumnya. Pembahasan mencakup rincian implementasi sistem yang meliputi pemodelan graf ekonomi, simulasi berbasis agen, dan integrasi algoritma evolusioner. Selanjutnya disajikan hasil pelaksanaan eksperimen pada model ekonomi generatif serta studi kasus menggunakan data ekonomi gim eksternal melalui \textit{game design document}, yang kemudian dianalisis untuk mengevaluasi efektivitas pendekatan yang diusulkan.

\section{Implementasi Model Sistem \textit{Reward} dan Lingkungan Simulasi}
Pada bagian ini, sistem reward dari ekonomi gim diimplementasikan menggunakan paradigma Pemrograman Berorientasi Objek (Object-Oriented Programming). Hubungan antar komponen dalam arsitektur program ini dirancang menggunakan prinsip pewarisan (inheritance) untuk memastikan modularitas. Struktur kelas tersebut secara visual direpresentasikan pada Gambar \ref{fig:arsitektur_kelas} (\nameref{fig:arsitektur_kelas}), yang memperlihatkan bagaimana kelas abstrak Node menjadi fondasi bagi seluruh entitas ekonomi dalam sistem. Sistem ekonomi dimodelkan sebagai Jaringan Logika Aliran Dinamika Skalar yang direpresentasikan melalui struktur data graf berarah (Directed Graph), di mana $G = (V, E, W)$ merupakan gabungan dari himpunan titik simpul (\textit{Nodes/Vertices}), himpunan sisi koneksi (Edges), dan matriks bobot (\textit{Weights}).

\subsection{Representasi Struktur Graf}
Struktur utama sistem dibangun di dalam kelas Graph yang bertugas sebagai pengelola seluruh daftar koneksi antar \textit{node} serta mengatur pola struktur jaringan tersebut. Kelas ini menginisialisasi arsitektur graf berdasarkan konfigurasi awal (config) yang mendefinisikan himpunan kuantitas tiap tipe titik. Dalam pembentukan graf, bobot atau nilai sisi dievaluasi menggunakan dua pendekatan:
\begin{enumerate}
    \item Bobot Deterministik: Nilai skalar riil yang merepresentasikan jumlah aliran sumber daya pasti yang dipindahkan antar simpul. Pada saat inisialisasi awal, nilai ini dapat dibangkitkan menggunakan distribusi seragam (uniform distribution).
    \item Koefisien Stokastik: Probabilitas khusus untuk simpul percabangan stokastik (rantai Markov). Untuk memastikan probabilitas transisi tetap valid, nilai probabilitas $p_i$ dibangkitkan menggunakan Parameter Distribusi Dirichlet (Distribusi Probabilitas Simplex n-dimensi) yang menjamin kriteria integral $\sum p_i = 1$ selalu terpenuhi.

\begin{figure}[htbp]
\begin{tikzpicture}[
  scale=1,
  transform shape,
  font=\sffamily\tiny,
  class/.style={
    draw=blue!50!black, thick,
    rectangle split, rectangle split parts=3,
    rectangle split part fill={blue!10, blue!5, white},
    text width=3.2cm, align=center, rounded corners=2pt
  },
  abstract/.style={
    draw=gray!70, thick, dashed,
    rectangle split, rectangle split parts=3,
    rectangle split part fill={gray!20, gray!10, white},
    text width=3.4cm, align=center, rounded corners=2pt
  },
  inherit/.style ={draw, thick,
    -{Triangle[open,fill=white,length=3.5mm,width=3mm]}},
  compose/.style ={draw, thick,
    -{Diamond[fill=black,length=3mm,width=2.2mm]}},
  associate/.style={draw, thick, -stealth}
]
 
%% ── 0a. AgentSimulator (extends Simulator) ────────────────────────
\node[class] (AgentSimulator) at (2.0,0) {
  \textbf{AgentSimulator}
  \nodepart{second}
  \begin{tabular}[t]{@{}l@{}}
    + agent : Agent \\
    + graph : Graph \\
    + monitoring : dict
  \end{tabular}
  \nodepart{third}
  \begin{tabular}[t]{@{}l@{}}
    + run(steps)
  \end{tabular}
};

%% ── 0b. Agent ─────────────────────────────────────────────────────
\node[class] (Agent) at (2.0,-3.8) {
  \textbf{Agent}
  \nodepart{second}
  \begin{tabular}[t]{@{}l@{}}
    + behavior : str
  \end{tabular}
  \nodepart{third}
  \begin{tabular}[t]{@{}l@{}}
    + play\_step(graph)
  \end{tabular}
};

%% ── 1. Simulator ──────────────────────────────────────────────────
\node[class] (Simulator) at (7.4,0) {
  \textbf{Simulator}
  \nodepart{second}
  \begin{tabular}[t]{@{}l@{}}
    + graph : Graph \\
    + monitoring : dict
  \end{tabular}
  \nodepart{third}
  \begin{tabular}[t]{@{}l@{}}
    + run(steps) \\
    + monitor() \\
    + plot\_monitor()
  \end{tabular}
};
 
%% ── 2. Graph ──────────────────────────────────────────────────────
\node[class] (Graph) at (7.4,-3.8) {
  \textbf{Graph}
  \nodepart{second}
  \begin{tabular}[t]{@{}l@{}}
    + config : dict \\
    + edge\_list : list \\
    + nodes : list \\
    + weights : list
  \end{tabular}
  \nodepart{third}
  \begin{tabular}[t]{@{}l@{}}
    + init\_nodes() \\
    + set\_edge\_weights() \\
    + simulate() \\
    + plot()
  \end{tabular}
};
 
%% ── 3. Node (Abstract) ────────────────────────────────────────────
\node[abstract] (Node) at (7.4,-8.8) {
  \textit{\textlangle\textlangle Abstract\textrangle\textrangle}\\\textbf{Node}
  \nodepart{second}
  \begin{tabular}[t]{@{}l@{}}
    + id : string \\
    + input\_edges : list \\
    + output\_edges : list
  \end{tabular}
  \nodepart{third}
  \begin{tabular}[t]{@{}l@{}}
    \textit{+ step(call\_chain)} \\
    + consume() \\
    + connect(node, value) \\
    + get\_state()
  \end{tabular}
};
 
%% ── 4. Edge ───────────────────────────────────────────────────────
\node[class] (Edge) at (2.4,-8.8) {
  \textbf{Edge}
  \nodepart{second}
  \begin{tabular}[t]{@{}l@{}}
    + value : float \\
    + node : Node
  \end{tabular}
  \nodepart{third}
  \phantom{x}
};
 
%% ── 5. Source ─────────────────────────────────────────────────────
\node[class] (Source) at (0.8,-14.0) {
  \textbf{Source}
  \nodepart{second}
  \phantom{x}
  \nodepart{third}
  \begin{tabular}[t]{@{}l@{}}
    + drop\_to(edge)
  \end{tabular}
};
 
%% ── 6. Pool ───────────────────────────────────────────────────────
\node[class] (Pool) at (4.8,-14.0) {
  \textbf{Pool}
  \nodepart{second}
  \begin{tabular}[t]{@{}l@{}}
    + pool : float
  \end{tabular}
  \nodepart{third}
  \begin{tabular}[t]{@{}l@{}}
    + consume(value) \\
    + reset()
  \end{tabular}
};
 
%% ── 7. Converter ──────────────────────────────────────────────────
\node[class] (Converter) at (9.4,-14.0) {
  \textbf{Converter}
  \nodepart{second}
  \begin{tabular}[t]{@{}l@{}}
    + is\_auto : bool \\
    + called : bool
  \end{tabular}
  \nodepart{third}
  \begin{tabular}[t]{@{}l@{}}
    + consume()
  \end{tabular}
};
 
%% ── 8. RandomGate ─────────────────────────────────────────────────
\node[class] (RandomGate) at (13.6,-14.0) {
  \textbf{RandomGate}
  \nodepart{second}
  \phantom{x}
  \nodepart{third}
  \begin{tabular}[t]{@{}l@{}}
    + consume(entity)
  \end{tabular}
};
 
%% ── 9. FixedPool ──────────────────────────────────────────────────
\node[class] (FixedPool) at (2.4,-18.8) {
  \textbf{FixedPool}
  \nodepart{second}
  \phantom{x}
  \nodepart{third}
  \begin{tabular}[t]{@{}l@{}}
    + get\_fix()
  \end{tabular}
};
 
%% ── 10. Drain ─────────────────────────────────────────────────────
\node[class] (Drain) at (7.8,-18.8) {
  \textbf{Drain}
  \nodepart{second}
  \phantom{x}
  \nodepart{third}
  \begin{tabular}[t]{@{}l@{}}
    + consume()
  \end{tabular}
};
 
%% ═══════════════ RELASI ═══════════════
 
%% AgentSimulator -▷ Simulator (pewarisan)
\draw[inherit] (AgentSimulator.east) -- (Simulator.west)
  node[midway,above=2pt,font=\tiny]{extends};

%% AgentSimulator → Agent (asosiasi kepemilikan)
\draw[associate] (AgentSimulator.south) -- (Agent.north)
  node[midway,right=2pt,font=\tiny]{menggunakan};

%% Simulator → Graph (asosiasi)
\draw[associate] (Simulator.south) -- (Graph.north)
  node[midway,right=2pt,font=\tiny]{mengoperasikan};
 
%% Graph ─◆ Node (komposisi 1..*)
\draw[compose] (Node.north) -- (Graph.south)
  node[midway,right=2pt,font=\tiny]{1..*};
 
%% Edge ─◆ Node (komposisi 0..*)
\draw[compose] (Edge.east) -- (Node.west)
  node[midway,above=2pt,font=\tiny]{0..*};
 
%% Edge.node → Node (asosiasi balik)
\draw[associate]
  (Edge.north) -- ++(0,1.5)
  -| ([xshift=-1.2cm]Node.north);
 
%% Pewarisan ke Node
\coordinate (jN) at (7.4,-11.8);
\coordinate (topLine) at (1,-11.9); % sesuaikan y

\draw[thick] (Source.north)     |- (topLine) -| (jN);
\draw[thick] (Pool.north)       |- (topLine) -| (jN);
\draw[thick] (Converter.north)  |- (topLine) -| (jN);
\draw[thick] (RandomGate.north) |- (topLine) -| (jN);
\draw[inherit] (jN) -- (Node.south);
 
%% Pewarisan ke Pool
\coordinate (jP) at (4.8,-17.1);
\draw[thick] (FixedPool.north) -- ++(0,0.6) -| (jP);
\draw[thick] (Drain.north)     -- ++(0,0.6) -| (jP);
\draw[inherit] (jP) -- (Pool.south);

\begin{scope}[shift={(10.8,0.0)}]
  % Kotak luar legenda
  \draw[draw=gray!60, rounded corners=3pt, fill=gray!5]
    (0,0) rectangle (5.8,-4.2);
 
  % Judul
  \node[font=\sffamily\small\bfseries, anchor=west] at (0.25,-0.45)
    {Legenda};
  \draw[gray!40] (0.2,-0.75) -- (5.6,-0.75);
 
  % ── Baris 1: Asosiasi ──────────────────────────────────────────
  \draw[associate] (0.3,-1.25) -- (1.6,-1.25);
  \node[font=\sffamily\tiny, anchor=west] at (1.75,-1.25)
    {Asosiasi (menggunakan)};
 
  % ── Baris 2: Komposisi ─────────────────────────────────────────
  \draw[compose] (0.3,-2.05) -- (1.6,-2.05);
  \node[font=\sffamily\tiny, anchor=west] at (1.75,-2.05)
    {Komposisi (memiliki)};
 
  % ── Baris 3: Pewarisan ─────────────────────────────────────────
  \draw[inherit] (0.3,-2.85) -- (1.6,-2.85);
  \node[font=\sffamily\tiny, anchor=west] at (1.75,-2.85)
    {Pewarisan / Inheritance};
 
  % ── Baris 7: Kelas abstrak ──────────────────────────────────────
  \draw[gray!70, thick, dashed, rounded corners=2pt]
    (0.3,-3.65) rectangle (1.6,-3.35);
  \node[font=\sffamily\tiny\itshape] at (0.95,-3.6)
    {};
  \node[font=\sffamily\tiny, anchor=west] at (1.75,-3.5)
    {Kelas abstrak};
 
\end{scope}
 
\end{tikzpicture}
\caption{Arsitektur Kelas Diagram}
\label{fig:arsitektur_kelas}
\end{figure}
    
\end{enumerate}
Graf ini juga dilengkapi dengan mekanisme normalisasi internal (scaling) untuk memverifikasi validitas properti ruang sampel probabilitas secara real-time setiap kali terjadi pembaruan atau mutasi bobot sisi selama proses simulasi berjalan.

\begin{figure}[htbp]
    \centering
    \includegraphics[width=0.7\textwidth]{tugas-akhir/assets/images/result_2.png}
    \caption{Contoh Visualisasi Struktur Graf Ekonomi Gim}
    \label{fig:initial_graph}
\end{figure}

\subsection{Implementasi \textit{Node} Ekonomi}
Setiap entitas ekonomi diwariskan (polimorfisme) dari abstraksi basis kelas \texttt{Node}. Kelas dasar ini mengatur konstrain domain fungsi input dan kodomain output, serta menerapkan hubungan antar entitas tersusun dengan benar. Evaluasi koneksi memuat Aksioma Irreflexive/Acyclic yang secara ketat menolak pembentukan loop tertutup (infinite recursion) pada simpul yang sama, sehingga arsitektur tetap berwujud \textit{Directed Acyclic Graph} (DAG).

Beberapa turunan simpul utama yang diimplementasikan meliputi:

\begin{itemize}
    \item \textit{Source},  bertindak sebagai batas titik awal ($S_0$) yang mensuplai aliran sumber daya ke dalam sistem dengan nilai in-degree sama dengan 0. Node ini mendorong transformasi nilai output maju ke subgraf selanjutnya.
    \item \textit{Pool}, simpul penyimpanan yang mengakumulasi sumber daya melalui operasi aljabar penjumlahan nilai state. Persamaan beda yang diterapkan adalah $\Delta P = P(t_{n+1}) = P(t_n) + x_{in} - x_{out}$. Terdapat variasi FixedPool yang mengimplementasikan batasan fungsi non-linier terikat $f(x) = \min(\text{pool} + x, \text{fix})$ untuk membatasi kapasitas maksimum inventaris agen.
    \item \textit{Converter}, mengeksekusi fungsi transformasi $C(x)$ yang mengubah kombinasi elemen input menjadi output baru. Simpul ini memberlakukan operator himpunan Boolean AND, di mana konversi hanya dieksekusi apabila seluruh prasyarat edge input memenuhi batas nilai $x \geq e$.
    \item \textit{RandomGate}, memodelkan percabangan aliran terbobot (weighted stochastic split). Distribusi aliran dieksekusi menggunakan percobaan probabilitas majemuk distribusi multinomial, yang mengukur vektor sampel observasi berdasarkan koefisien proporsi kejadian $P \in [0, 1]$.
    \item \textit{Drain}, Simpul sink dengan nilai out-degree sama dengan 0 yang merepresentasikan pengeluaran tak terhingga atau pembuangan variabel dari sistem persamaan diferensial.
\end{itemize}

\subsection{Implementasi Mekanisme Aliran (\textit{Edge}) dan Simulasi Waktu}
Dinamika aliran dari waktu ke waktu dieksekusi oleh mesin komputasi numerik diskret dalam kelas \texttt{Simulator} diperluas (inheritance) menjadi kelas \texttt{AgentSimulator}. Simulator ini menjalankan perulangan untuk mengevaluasi pergeseran dimensi waktu dari $t = 0$ hingga $t = n$ langkah.
% , yang dimodelkan sebagai sistem persamaan beda orde pertama:
% \[
% \mathbf{x}_{t+1} = f(\mathbf{x}_t), \quad t = 0,1,2,\dots,n-1.
% \]

Mekanisme relasi transmisi ditangani oleh objek \texttt{Edge} yang meneruskan vektor besaran dari titik $u$ ke titik tujuan $v$. 

\begin{algorithm}[h]
\caption{Simulasi Dinamika Aliran Graf Waktu Diskret}
\label{alg:graph-simulator}
\begin{algorithmic}[1]
\INPUT Graf $G = (V, E)$, jumlah langkah $n$ 
\OUTPUT Deret waktu $\mathbf{X}(t)$ untuk setiap node \textit{Pool}

\State Ekstrak himpunan node: $P$ (\textit{Pool}), $S$ (\textit{Source}), $C$ (\textit{Converter})
\State Inisialisasi struktur monitoring $\mathbf{X} \gets \emptyset$

\For{$t = 0$ hingga $n-1$}
    \For{setiap $s \in S$}
        \State $s.\text{step}()$ \Comment{Injeksi input dan propagasi dalam graf}
    \EndFor
    
    \For{setiap $p \in P$}
        \State Rekam nilai state $p$ ke $\mathbf{X}(t)$
    \EndFor
    
    \For{setiap $c \in C$}
        \State $c.\text{called} \gets \text{False}$ \Comment{Reset status evaluasi}
    \EndFor
\EndFor

\For{setiap $p \in P$}
    \State Reset state $p$
\EndFor

\State \Return $\mathbf{X}$
\end{algorithmic}
\end{algorithm}

Pada setiap iterasi waktu $t$:

\begin{enumerate}
    \item Injeksi Batas, Variabel bebas batas (\textit{boundary parameters}) dimasukkan melalui simpul \texttt{Source}.
    \item Propagasi (\textit{Forward Pass}), Nilai ditransmisikan melintasi sisi graf, mengubah state inventaris pada \texttt{Pool} dan memicu evaluasi transformasi pada \texttt{Converter} jika ketersediaan sumber daya tercukupi.
    \item Perekaman (Monitoring), Setelah propagasi selesai pada satu langkah waktu, simulator mengeksekusi metode pencatatan (snapshot) untuk merekam metrik $P(t)$ dari setiap simpul penyimpanan. Array historis ini membentuk vektor seri deret waktu $\vec{X}(t)$.
    \item Reset State, Flag status pemanggilan untuk setiap simpul perantara diatur ulang menjadi \textit{False} untuk mempersiapkan evaluasi transformasi pada interval $t+1$.
\end{enumerate}

Kumpulan matriks historis dari hasil relasi transisi ini kemudian diekspor untuk diproyeksikan sebagai grafik deret waktu, yang memvisualisasikan bagaimana aliran ekonomi berevolusi secara sekuensial.

\subsection{Implementasi  Simulasi Perilaku dan Entitas Agen}

Simulasi tidak bekerja secara pasif, melainkan dikendalikan oleh entitas otonom yang memodelkan pemain gim sungguhan menggunakan kelas \texttt{Agent}. Agen ini bertugas mengawasi simpul \texttt{Converter} yang membutuhkan pemicu manual dan tidak digerakkan murni oleh probabilitas gerbang Markov (\texttt{RandomGate}).

Logika pengambilan keputusan agen diklasifikasikan ke dalam tiga profil perilaku yang memengaruhi dinamika transisi state pada graf:

\begin{itemize}
    \item Agresif, pada model ini agen akan langsung memicu semua fungsi konversi secara maksimum pada setiap langkah waktu $t$, asalkan batas bawah prasyarat sumber daya telah terpenuhi. Model ini merepresentasikan pemain yang mengeksekusi aksi secepat mungkin tanpa menunda.
    \item Pasif, model ini mengevaluasi ketersediaan sumber daya menggunakan operasi logika Boolean AND terhadap seluruh simpul konversi manual. Agen hanya akan mengeksekusi aksi jika semua opsi konversi (\textit{skill} atau opsi \textit{crafting}) telah siap (sumber daya mencukupi dan \textit{cooldown} selesai). Jika proposisi bernilai \textit{True}, seluruh konversi dieksekusi secara bersamaan. Profil ini merepresentasikan pemain yang bermain aman dan menunggu momen optimal.
    \item Probabilistik, Pemodelan ini menggunakan pendekatan berbasis Monte-Carlo. Untuk setiap simpul konversi manual yang telah memenuhi syarat sumber daya, agen mengeksekusinya berdasarkan probabilitas sukses sebesar $p = 0.5$ (menggunakan distribusi seragam). Profil ini menyimulasikan pemain kasual yang pengambilan keputusannya tidak menentu dan probabilistik.
\end{itemize}

Integrasi antara \texttt{AgentSimulator} dan \texttt{Agent} memungkinkan algoritma optimasi nantinya tidak hanya menguji ekonomi yang berjalan secara matematis murni, tetapi juga menguji apakah ekonomi gim tersebut stabil terhadap berbagai variasi gaya bermain pemain.

\paragraph{Contoh Trace Simulasi Agen.}
Sebagai ilustrasi perbedaan perilaku antar profil agen, perhatikan contoh graf sederhana berikut. Graf memiliki satu \textit{Source} yang menginjeksikan 10 unit per langkah ke \textit{Pool\_A}. Terdapat dua \textit{Converter} manual yang keduanya menarik dari \textit{Pool\_A}: \textit{C1} membutuhkan 8 unit dan menghasilkan 12 unit ke \textit{Pool\_Out1}, serta \textit{C2} membutuhkan 10 unit dan menghasilkan 8 unit ke \textit{Pool\_Out2}. Tabel \ref{tab:agent_trace} menunjukkan dinamika sistem untuk tiga profil agen selama tiga langkah waktu.

\begin{table}[htbp]
\centering
\caption{Trace Simulasi Tiga Profil Agen pada Graf Contoh ($x = 30$)}
\label{tab:agent_trace}
\begin{adjustbox}{max width=\textwidth}
\begin{tabular}{clcccccc}
\hline
$t$ & \textbf{Profil} & \textbf{Pool\_A} & \textbf{C1 aktif?} & \textbf{C2 aktif?} & \textbf{Pool\_A akhir} & \textbf{Pool\_Out1} & \textbf{Pool\_Out2} \\
\hline
\multirow{3}{*}{0} & Aggressive & 10 & Ya (10$\geq$8)  & Tidak (10$<$10) & 2  & 12 & 0  \\
                   & Passive    & 10 & \multicolumn{2}{c}{Tidak (C2 belum siap)} & 10 & 0  & 0  \\
                   & Random     & 10 & \multicolumn{2}{c}{Masing-masing $p=0.5$} & —  & —  & —  \\
\hline
\multirow{3}{*}{1} & Aggressive & 2+10=12  & Ya (12$\geq$8) & Tidak (12$<$10) & 4  & 24 & 0  \\
                   & Passive    & 10+10=20 & \multicolumn{2}{c}{Tidak (C2 belum siap, 20$\geq$10?)} & — & — & — \\
                   & Random     & —        & \multicolumn{2}{c}{Masing-masing $p=0.5$} & — & — & — \\
\hline
\multirow{2}{*}{2} & Aggressive & 4+10=14  & Ya (14$\geq$8) & Ya (14$\geq$10) & $14{-}8{-}10=\!-4$* & 36 & 8 \\
                   & Passive    & 20+10=30 & \multicolumn{2}{c}{Ya (C1 dan C2 keduanya siap)} & $30{-}8{-}10=12$ & 12 & 8 \\
\hline
\end{tabular}
\end{adjustbox}
\end{table}

\textit{*Catatan: Pada kondisi riil, Balancer mengoptimasi bobot sehingga kasus pool negatif dihindari.}

Dari trace di atas terlihat perbedaan mendasar antar profil. Agen \textit{aggressive} langsung memicu C1 di $t=0$ begitu syaratnya terpenuhi, menghasilkan akumulasi \textit{Pool\_Out1} yang lebih cepat — namun dapat menyebabkan pool terkuras sebelum konverter lain siap. Agen \textit{passive} menunggu hingga \textit{semua} konverter siap secara bersamaan (baru $t=2$), sehingga menghasilkan output yang lebih teratur namun lebih lambat. Perbedaan pola interaksi inilah yang menciptakan variasi sinyal fitness yang diterima oleh algoritma evolusioner, mendorong eksplorasi ruang parameter yang lebih beragam dibandingkan simulasi tanpa agen.


% \section{Implementasi Simulasi Agen}
% Bagian implementasi ini menjabarkan bagaimana entitas agen dikonstruksi untuk berinteraksi di dalam lingkungan graf ekonomi. Agen, direpresentasikan dalam modul \texttt{geevo.agent\_simulation}, berinteraksi dengan tipe perilaku berbeda.

% \subsection{Pembuatan Entitas Agen}
% Kelas \texttt{Agent} menampung atribut kekayaan internal pemain (\textit{wealth}/\textit{inventory}) dan parameter preferensi yang menetapkan kecenderungannya terhadap risiko dan eksploitasi fitur dari fungsi ekonomi. Agen diinisiasi dalam variasi profil profil pemain tertentu seperti profil \textit{aggressive}, \textit{passive}, dan \textit{random}.

% \subsection{Logika Interaksi dan Kondisi \textit{State}}
% Dinamika \textit{loop} agen mensimulasikan akumulasi perolehan sumber daya pemain pada setiap langkah waktu (\textit{step}). Pada simulasi ini, \textit{auto-farming} dikontrol oleh \textit{flag} internal sementara aksi penukaran dikalkulasikan sesuai probabilitas preferensi pemain.

\section{Implementasi dan Integrasi Algoritma Evolusioner}
Algoritma evolusioner digunakan untuk secara berulang menemukan konfigurasi laju terbaik yang menyeimbangkan sistem. Komponen evolusi dienkapsulasi pada kelas \texttt{Balancer}.

\subsection{Representasi Kromosom Konfigurasi \textit{Reward}}
Vektor numerik merepresentasikan bobot probabilitas, laju \textit{source}, serta harga komoditas (laju \textit{sink}/\textit{drain}). Vektor ini bermutasi dan disilangkan pada proses reproduksi algoritma penyeimbangan.

\subsection{Fungsi \textit{Fitness}}
Nilai \textit{fitness} dievaluasi dari performa ekonomi melalui simulasi singkat dengan kondisi terukur. Keseimbangan (fitness tertinggi) didefinisikan sebagai tingkat stabilitas terbaik di mana sistem terhindar dari inflasi maupun krisis sumber daya ekstrem.

Dua fungsi fitness diimplementasikan dalam penelitian ini untuk membandingkan efektivitas pendekatan evaluasi yang berbeda:

\begin{enumerate}
    \item \textbf{Fungsi Objektif 3 (obj3)}, merupakan fungsi proporsional yang diadaptasi dari kerangka GEEvo \citep{Rupp_2024}. Nilai fitness diukur berdasarkan rasio proporsional antara nilai aktual pool $s_t$ pada langkah waktu akhir dan nilai target $x$:
    \[
        f_3(s_t, x) = \frac{\min(s_t,\; x)}{\max(s_t,\; x)} \longrightarrow \max \approx 1.0
    \]
    Sistem dinyatakan seimbang (\textit{balanced}) apabila nilai fitness mencapai ambang $f_3 \geq 1 - \alpha$.

    \item \textbf{Fungsi Objektif 4 (obj4)}, merupakan kontribusi penelitian ini sebagai varian obj3 dengan mekanisme \textit{tolerance-threshold}. Fungsi ini memberikan nilai fitness sempurna ($f = 1.0$) apabila nilai aktual berada dalam toleransi $\alpha$ terhadap target, dan kembali ke formula proporsional apabila di luar batas tersebut:
    \[
        f_4(s_t, x) =
        \begin{cases}
            1.0 & \text{jika } |x - s_t| \leq \alpha \cdot x \\[4pt]
            \dfrac{\min(s_t,\; x)}{\max(s_t,\; x)} & \text{sebaliknya}
        \end{cases}
    \]
    Modifikasi ini menciptakan \textit{reward plateau} di sekitar nilai target, sehingga memberikan gradien sinyal yang lebih informatif bagi algoritma evolusioner untuk konvergen lebih cepat.
\end{enumerate}

Fitness dievaluasi melalui $n_{\text{sim}} = 10$ simulasi Monte-Carlo per individu untuk mereduksi bias variansi akibat perilaku agen yang stokastik. Nilai fitness akhir merupakan rata-rata dari seluruh simulasi tersebut.

Sebagai ilustrasi perbedaan konkret antara obj3 dan obj4, perhatikan contoh berikut. Misalkan target $x = 30$, $\alpha = 0.05$, sehingga toleransi obj4 adalah $\alpha \cdot x = 0.05 \times 30 = 1.5$ dan ambang keseimbangan obj3 adalah $1 - \alpha = 0.95$.

\begin{table}[htbp]
\centering
\caption{Perbandingan Nilai obj3 dan obj4 untuk Berbagai $s_t$ ($x = 30,\ \alpha = 0.05$)}
\label{tab:obj_example}
\begin{tabular}{ccccc}
\hline
\textbf{Nilai Pool} $s_t$ & \textbf{$f_3$} & \textbf{$f_4$} & \textbf{Seimbang (obj3, $\geq 0.95$)} & \textbf{Seimbang (obj4, $|30{-}s_t| \leq 1.5$)} \\
\hline
10.0 & 0.333 & 0.333 & Tidak & Tidak \\
25.0 & 0.833 & 0.833 & Tidak & Tidak \\
28.0 & 0.933 & 0.933 & Tidak & Tidak \\
28.5 & 0.950 & \textbf{1.000} & Ya    & Ya    \\
29.0 & 0.967 & \textbf{1.000} & Ya    & Ya    \\
30.0 & 1.000 & \textbf{1.000} & Ya    & Ya    \\
31.0 & 0.968 & \textbf{1.000} & Ya    & Ya    \\
31.5 & 0.952 & \textbf{1.000} & Ya    & Ya    \\
32.0 & 0.938 & 0.938 & Tidak & Tidak \\
35.0 & 0.857 & 0.857 & Tidak & Tidak \\
\hline
\end{tabular}
\end{table}

Tabel \ref{tab:obj_example} menunjukkan perbedaan kritis antara kedua fungsi. Untuk $s_t = 29$ misalnya, obj3 menghasilkan fitness $29/30 = 0.967$, sementara obj4 memberikan $1.000$ karena $|30 - 29| = 1 \leq 1.5$. Keduanya menyatakan kondisi \textit{balanced}, tetapi sinyal yang dikirimkan berbeda: obj4 memberikan sinyal tegas ``sudah optimal'' sehingga algoritma dapat berhenti lebih awal. Sebaliknya, untuk $s_t = 32$ di mana pool \textit{melebihi} target, obj3 memberi 0.938 (tidak seimbang) dan obj4 juga memberi 0.938 karena $|30 - 32| = 2 > 1.5$. Zona keseimbangan obj4 bersifat simetris terhadap target (28.5 hingga 31.5), memberikan tekanan seleksi yang setara untuk pool yang kekurangan maupun kelebihan sumber daya.

Kurva pencarian penyeimbangan model utama ditunjukkan pada pencarian \texttt{EvoGraphGenerator} di Gambar \ref{fig:egg_fitness}.

\begin{figure}[htbp]
    \centering
    \includegraphics[width=0.6\textwidth]{tugas-akhir/assets/images/result_1.png}
    \caption{Historis Evaluasi \textit{Fitness} Algoritma Evolusi Generatif}
    \label{fig:egg_fitness}
\end{figure}

\section{Pelaksanaan Uji Coba dan Studi Kasus}
% Pelaksanaan dilakukan dengan menyimulasikan iterasi pada beberapa profil perilaku agen dan dengan desain graf sebelum maupun sesudah penerapan optimasi algoritma.

\subsection{Deskripsi Data \textit{Game Design Document}}

Pengujian pada data eksternal dilakukan menggunakan dua \textit{Game Design Document} (GDD) nyata yang dipilih berdasarkan perbedaan kompleksitas sistem ekonominya. Pemilihan dua GDD yang kontras bertujuan untuk menguji kemampuan generalisasi model terhadap topologi ekonomi yang beragam.

\begin{enumerate}
    \item \textbf{MMORPG Game Design Document} (\texttt{MMORPG\_GameDesign\_v27.pdf}), merupakan dokumen desain untuk gim \textit{massively multiplayer online role-playing game} dengan sistem ekonomi kompleks yang mencakup mekanisme pengumpulan mata uang emas, penyimpanan sumber daya bertingkat, \textit{crafting}, \textit{upgrading}, dan distribusi probabilistik hadiah. Proses ekstraksi kata kunci menghasilkan konfigurasi node $\{$\textit{Source}: 2, \textit{RandomGate}: 3, \textit{Pool}: 5, \textit{Converter}: 3$\}$, menghasilkan total 13 node yang merepresentasikan kompleksitas ekonomi kelas MMORPG.

    \item \textbf{VHS Horror Game Design Document} (\texttt{VHS-Horror-Game-Design-Document.pdf}), merupakan dokumen desain untuk gim \textit{survival horror} dengan sistem ekonomi yang lebih sederhana, berfokus pada inventori terbatas dan mekanisme bertahan hidup. Ekstraksi menghasilkan konfigurasi node $\{$\textit{Source}: 1, \textit{RandomGate}: 2, \textit{Pool}: 3, \textit{Converter}: 1$\}$, menghasilkan total 7 node yang mencerminkan minimnya sumber daya dalam gim bertipe \textit{survival}.
\end{enumerate}

Untuk setiap dataset GDD, sebanyak $n = 10$ graf ekonomi dibangkitkan menggunakan \texttt{EvolutionaryGraphGeneration} (EGG) berdasarkan konfigurasi yang diekstrak. Seluruh eksperimen pada dataset GDD menggunakan kondisi TA penuh (Kondisi C: simulasi agen + obj4) sebagai representasi kontribusi akhir penelitian ini.

\subsection{Mekanisme Penerjemahan GDD Menjadi Bentuk Graf}

Penerjemahan GDD menjadi konfigurasi graf ekonomi dilakukan oleh modul \texttt{GDDExtractor} yang dikembangkan sebagai kontribusi penelitian ini. Proses penerjemahan berlangsung dalam tiga tahap:

\begin{enumerate}
    \item \textbf{Ekstraksi Teks PDF}, modul membaca teks seluruh halaman dokumen GDD menggunakan pustaka \textit{PyMuPDF} atau \textit{PyPDF2} sebagai alternatif. Teks seluruh dokumen diproses dalam bentuk huruf kecil (\textit{lowercase}) untuk normalisasi pencocokan kata kunci.

    \item \textbf{Pemetaan Kata Kunci (\textit{Keyword Matching})}, setiap tipe node GEEvo dipetakan ke satu himpunan kata kunci yang merepresentasikan konsep ekonomi relevan dalam bahasa natural. Frekuensi kemunculan kata kunci di seluruh dokumen dihitung untuk setiap tipe node, sebagaimana dirinci pada Tabel \ref{tab:keyword_map}.

    \item \textbf{Konversi Proporsional ke Konfigurasi Node}, hasil perhitungan frekuensi dikonversi secara proporsional menjadi jumlah node setiap tipe. Setiap tipe node dijamin mendapat minimal satu node, dan total node dibatasi sesuai parameter \texttt{max\_nodes}. Konfigurasi yang dihasilkan selanjutnya digunakan sebagai masukan EGG untuk membangkitkan topologi graf.
\end{enumerate}

\begin{table}[htbp]
\centering
\caption{Pemetaan Kata Kunci ke Tipe Node GEEvo}
\label{tab:keyword_map}
\begin{tabular}{ll}
\hline
\textbf{Tipe Node} & \textbf{Contoh Kata Kunci} \\
\hline
\textit{Source}     & drop, reward, quest, loot, gather, generate, spawn, earn, produce \\
\textit{Pool}       & gold, currency, resource, inventory, wallet, balance, mana, xp, health \\
\textit{Converter}  & shop, crafting, trade, market, upgrade, buy, sell, convert, forge \\
\textit{RandomGate} & chance, luck, rng, random, gamble, gacha, dice, drop rate, loot table \\
\hline
\end{tabular}
\end{table}

\begin{algorithm}[h]
\caption{Ekstraksi Konfigurasi GEEvo dari GDD}
\label{alg:gdd-extractor}
\begin{algorithmic}[1]
\INPUT Jalur file GDD $P_{pdf}$, jumlah maksimum node $N_{max}$
\OUTPUT Konfigurasi node $C = \{$\textit{Source}, \textit{Pool}, \textit{Converter}, \textit{RandomGate}$\}$

\State Baca teks seluruh halaman GDD: $T \gets \text{read\_pdf}(P_{pdf})$
\For{setiap tipe node $k \in \{$\textit{Source, Pool, Converter, RandomGate}$\}$}
    \State $w_k \gets \sum_{\text{keyword} \in K_k} \text{count}(\text{keyword},\; T)$
\EndFor
\If{$\sum_k w_k = 0$}
    \State \Return konfigurasi \textit{fallback} default $\{$S:2, P:4, C:2, RG:1$\}$
\EndIf
\State Distribusikan $N_{max}$ node secara proporsional berdasarkan $w_k$
\State Jamin minimal 1 node per tipe
\State \Return konfigurasi $C$
\end{algorithmic}
\end{algorithm}

\subsection{Skenario Uji Coba Dinamika Graf Generatif}

Uji coba dilakukan dalam dua skenario utama yang bertujuan mengevaluasi efektivitas pendekatan dari sudut pandang berbeda:

\begin{enumerate}
    \item \textbf{Skenario Generatif (Abstrak)}, sebanyak $n = 30$ graf ekonomi dibangkitkan secara generatif menggunakan konfigurasi abstrak ($\textit{Source}: 3, \textit{RandomGate}: 2, \textit{Pool}: 4, \textit{Converter}: 1$). Skenario ini mengevaluasi kemampuan dasar sistem penyeimbangan pada topologi yang dikontrol. Sebagai pelengkap, dilakukan juga perbandingan tiga kondisi evaluasi (A, B, C) pada $n = 194$ graf abstrak dengan konfigurasi yang sama untuk mengisolasi kontribusi masing-masing komponen yang diusulkan.

    \item \textbf{Skenario GDD (Data Nyata)}, menggunakan dua GDD nyata (MMORPG dan VHS) sebagai masukan eksternal. Masing-masing menghasilkan $n = 10$ graf yang topologinya diekstrak otomatis dari konten dokumen, kemudian diseimbangkan menggunakan kondisi penelitian penuh (Kondisi C).
\end{enumerate}

Tiga kondisi evaluasi yang dibandingkan dalam penelitian ini adalah:
\begin{itemize}
    \item \textbf{Kondisi A} : tanpa agen (simulasi pasif), fungsi fitness obj3 (baseline replikasi paper)
    \item \textbf{Kondisi B} : dengan agen (\textit{aggressive}, \textit{passive}, \textit{random}), fungsi fitness obj3 (isolasi dampak agen)
    \item \textbf{Kondisi C} : dengan agen, fungsi fitness obj4 (kondisi penelitian penuh, kontribusi utama)
\end{itemize}

Gambar \ref{fig:unbalanced} dan Gambar \ref{fig:balanced_graph} menunjukkan perbandingan visual struktur aliran sumber daya sebelum dan sesudah proses penyeimbangan pada salah satu contoh graf generatif.

\begin{figure}[htbp]
    \centering
    \includegraphics[width=0.7\textwidth]{tugas-akhir/assets/images/result_3.png}
    \caption{Simulasi Interaksi Agen pada Kondisi Graf Belum Seimbang}
    \label{fig:unbalanced}
\end{figure}

\begin{figure}[htbp]
    \centering
    \includegraphics[width=0.7\textwidth]{tugas-akhir/assets/images/result_5.png}
    \caption{Visualisasi Graf Ekonomi Generatif dengan Bobot \textit{Reward} Optimal}
    \label{fig:balanced_graph}
\end{figure}

\section{Analisis Hasil}
Analisis ini membandingkan fluktuasi pasokan sistem serta evaluasi \textit{fitness} algoritma secara holistik dari model \textit{baseline} menuju model yang telah dioptimasi algoritma \textit{reward balancing}.

\subsection{Kestabilan Sumber Daya pada Ekonomi Gim}
Memperhatikan performa awal dari Gambar \ref{fig:unbalanced}, perbandingan stabilitas diuji simulasi pada graf berbobot optimal. Gambar \ref{fig:balanced_sim} menunjukkan bahwa model telah berhasil mencegah akumulasi eksponensial maupun penipisan sumber daya selama iterasi agen, dan menjaga ekonomi di kondisi stasioner dinamis.

\begin{figure}[htbp]
    \centering
    \includegraphics[width=0.7\textwidth]{tugas-akhir/assets/images/result_6.png}
    \caption{Kestabilan Dinamika Volume Pool setelah Optimasi Konfigurasi Harga}
    \label{fig:balanced_sim}
\end{figure}

\subsection{Analisis Dinamika dan Perilaku Agen}
Uji coba menggunakan agen \textit{aggressive}, \textit{passive}, dan \textit{random} memvalidasi bahwa sistem teroptimasi memberikan batas kewajaran ekonomi lintas jenis pemain, terhimpun pada Gambar \ref{fig:agent_behaviors}. Perbedaan kecuraman aliran sumber daya yang dapat dicapai menunjukkan kerentanan spesifik tiap tipe dapat dieksplorasi pemain pada area \textit{Converter} \textit{manual} tanpa mendisrupsi ekuilibrium seluruh ekosistem gim.

\begin{figure}[htbp]
    \centering
    \begin{minipage}{0.48\textwidth}
        \centering
        \includegraphics[width=\textwidth]{tugas-akhir/assets/images/result_7.png}
        \subcaption{Simulasi Agen \textit{Aggressive}}
    \end{minipage}
    \begin{minipage}{0.48\textwidth}
        \centering
        \includegraphics[width=\textwidth]{tugas-akhir/assets/images/result_8.png}
        \subcaption{Simulasi Agen \textit{Passive}}
    \end{minipage}
    \caption{Perbandingan Dinamika Interaksi Profil Agen Berbeda \textit{Post-Balance}}
    \label{fig:agent_behaviors}
\end{figure}

\subsection{Evaluasi Kinerja Algoritma Evolusioner (Konvergensi \textit{Fitness})}
Algoritma utama (\texttt{Balancer}) terbukti berhasil menemukan titik batas maksimum (penurunan beban deviasi hingga konvergen) di mana bobot \textit{sink} dan probabilitas dapat beradaptasi secara heuristik. Hal ini tergambar jelas pada kenaikan performa historis konvergensi populasi parameter \textit{balancer} dari simulasi menuju optimasi maksimal (Gambar \ref{fig:balancer_monitor}).

\begin{figure}[htbp]
    \centering
    \includegraphics[width=0.6\textwidth]{tugas-akhir/assets/images/result_4.png}
    \caption{Kurva Konvergensi Metrik Evaluasi Optimasi \textit{Balancer}}
    \label{fig:balancer_monitor}
\end{figure}

\subsection{Evaluasi Komparatif Tiga Kondisi Eksperimen}
\label{subsec:abc_comparison}

Untuk mengisolasi kontribusi masing-masing komponen yang diusulkan, dilakukan evaluasi komparatif terhadap tiga kondisi pada $n = 194$ graf abstrak dengan tiga nilai parameter $\alpha$. Hasil perbandingan lengkap ditunjukkan pada Tabel \ref{tab:abc_full}.

\begin{table}[htbp]
\centering
\caption{Hasil Lengkap Ketiga Kondisi Evaluasi ($n=194$ Graf Abstrak)}
\label{tab:abc_full}
\begin{tabular}{llccccc}
\hline
\textbf{Kondisi} & $\alpha$ & \textbf{Bal.\%} & \textbf{Init.Bal.\%} & \textbf{Impr.\%} & \textbf{Std Dev} & \textbf{Med.Gen} \\
\hline
A: Tanpa Agen + obj3 & 0.05 & 48.97 & 3.09 & 78.35 & 0.293 & 500 \\
A: Tanpa Agen + obj3 & 0.01 & 14.95 & 1.03 & 82.99 & 0.245 & 500 \\
A: Tanpa Agen + obj3 & 0.00 &  4.64 & 0.52 & 80.41 & 0.231 & 500 \\
\hline
B: Agen + obj3       & 0.05 & 46.39 & 3.09 & 81.44 & 0.243 & 500 \\
B: Agen + obj3       & 0.01 & 12.37 & 1.55 & 86.60 & 0.192 & 500 \\
B: Agen + obj3       & 0.00 &  3.09 & 0.00 & 84.54 & 0.230 & 500 \\
\hline
\textbf{C: Agen + obj4 (TA)} & \textbf{0.05} & \textbf{54.64} & \textbf{3.09} & \textbf{87.11} & \textbf{0.240} & \textbf{8} \\
C: Agen + obj4 (TA)  & 0.01 & 15.46 & 0.52 & 84.54 & 0.235 & 500 \\
C: Agen + obj4 (TA)  & 0.00 &  3.09 & 0.00 & 84.54 & 0.230 & 500 \\
\hline
\end{tabular}
\end{table}

Tabel \ref{tab:abc_full} mengungkap dua temuan kunci tentang kontribusi masing-masing komponen:

\paragraph{Dampak Simulasi Agen (Kondisi B vs. A).}
Hasil pengujian menunjukkan bahwa Kondisi B secara konsisten mengungguli Kondisi A pada metrik \textit{Improved \%} di seluruh nilai $\alpha$. Peningkatan performa tercatat sebesar $3{,}09\%$ pada $\alpha=0{,}05$, $3{,}61\%$ pada $\alpha=0{,}01$, dan mencapai $4{,}13\%$ pada $\alpha=0{,}0$. Selain itu, nilai \textit{Fitness Std Dev} pada Kondisi B senantiasa lebih rendah dibandingkan Kondisi A, dengan perbedaan paling signifikan terlihat pada $\alpha=0{,}01$ ($0{,}192$ vs. $0{,}245$). Hal ini mengindikasikan bahwa penggunaan agen menghasilkan proses pencarian yang lebih stabil dan konsisten di berbagai topologi graf. Kontribusi agen ini terbukti bersifat robust dan stabil, terlepas dari tingkat keketatan ambang batas (\textit{threshold}) yang diterapkan.

\paragraph{Dampak \textit{obj4} (Kondisi C vs. B).}
Berbeda dengan simulasi agen, keunggulan fungsi \textit{obj4} ditemukan bersifat kondisional terhadap nilai $\alpha$. Keuntungan penggunaan \textit{obj4} terkonsentrasi pada $\alpha = 0{,}05$, yang ditandai dengan kenaikan \textit{Balanced \%} sebesar $8{,}25\%$, \textit{Improved \%} sebesar $5{,}67\%$, serta penurunan \textit{Median Generation} yang drastis dari $500$ menjadi $8$ generasi. Namun, pada $\alpha = 0{,}01$ dan $\alpha = 0{,}0$, hasil antara Kondisi B dan C cenderung identik. Fenomena ini terjadi karena rentang toleransi \textit{obj4} menjadi sangat sempit ($0{,}01 \times 30 = 0{,}3$ unit), sehingga mekanisme \textit{plateau} tidak lagi teraktivasi dan fungsi \textit{obj4} berperilaku ekuivalen dengan \textit{obj3}. Pola tersebut menunjukkan bahwa \textit{obj4} memberikan manfaat nyata pada \textit{threshold} yang lebih longgar, sementara pada \textit{threshold} ketat, gradien seleksi dari kedua fungsi tersebut menjadi serupa.

Gambar \ref{fig:convergence_abc} menunjukkan perbandingan kurva konvergensi dari ketiga kondisi tersebut.

\begin{figure}[htbp]
    \centering
    \includegraphics[width=0.85\textwidth]{tugas-akhir/assets/images/convergence_abc.png}
    \caption{Perbandingan Kurva Konvergensi \textit{Fitness} Ketiga Kondisi Evaluasi ($n=194$, $\alpha=0.05$)}
    \label{fig:convergence_abc}
\end{figure}

\subsection{Evaluasi Penyeimbangan pada Dataset GDD}
\label{subsec:gdd_results}

Evaluasi pada dataset GDD nyata dilakukan menggunakan kondisi TA penuh (Kondisi C). Tabel \ref{tab:mmorpg_results} dan Tabel \ref{tab:vhs_results} merangkum hasil penyeimbangan pada kedua dataset di berbagai nilai $\alpha$.

\begin{table}[htbp]
\centering
\caption{Hasil Evaluasi pada Dataset MMORPG GDD ($n=10$, Kondisi C)}
\label{tab:mmorpg_results}
\begin{tabular}{lccccc}
\hline
$\alpha$ & \textbf{Balanced \%} & \textbf{Initial Bal. \%} & \textbf{Improved \%} & \textbf{Median Gen} & \textbf{Exec Time (s)} \\
\hline
0.05 & 50.0\%  & 10.0\% & 70.0\%  & 251.0 & 3.609  \\
0.01 & 0.0\%   & 0.0\%  & 60.0\%  & 500.0 & 40.778 \\
0.0  & 10.0\%  & 0.0\%  & 80.0\%  & 500.0 & 54.733 \\
\hline
\end{tabular}
\end{table}

\begin{table}[htbp]
\centering
\caption{Hasil Evaluasi pada Dataset VHS Horror GDD ($n=10$, Kondisi C)}
\label{tab:vhs_results}
\begin{tabular}{lccccc}
\hline
$\alpha$ & \textbf{Balanced \%} & \textbf{Initial Bal. \%} & \textbf{Improved \%} & \textbf{Median Gen} & \textbf{Exec Time (s)} \\
\hline
0.05 & 30.0\% & 0.0\% & 100.0\% & 500.0 & 18.004 \\
0.01 & 20.0\% & 0.0\% & 100.0\% & 500.0 & 3.575  \\
0.0  & 0.0\%  & 0.0\% & 90.0\%  & 500.0 & 5.084  \\
\hline
\end{tabular}
\end{table}

\textbf{Dataset MMORPG.} Pada $\alpha = 0.05$, setengah dari total graf (50\%) berhasil mencapai kondisi seimbang, namun dengan median generasi 251 — jauh lebih lambat dibanding dataset abstrak (2.5 generasi). Perbedaan ini mencerminkan kompleksitas topologi MMORPG yang lebih tinggi (13 node, 3 \textit{Converter}) yang memperluas ruang pencarian algoritma. Sebesar 70\% graf tetap mengalami perbaikan (\textit{Improved}), menunjukkan bahwa algoritma efektif menemukan konfigurasi yang lebih baik meskipun tidak selalu mencapai ambang keseimbangan penuh. Kurva konvergensi dataset MMORPG ditunjukkan pada Gambar \ref{fig:konvergensi_mmorpg}.

\begin{figure}[htbp]
    \centering
    \includegraphics[width=0.75\textwidth]{tugas-akhir/assets/images/plot_konvergensi_mmorpg.png}
    \caption{Kurva Konvergensi \textit{Fitness} pada Dataset MMORPG GDD}
    \label{fig:konvergensi_mmorpg}
\end{figure}

\textbf{Dataset VHS Horror.} Dataset VHS menunjukkan pola yang berbeda: meskipun \textit{Balanced \%} relatif rendah (30\% pada $\alpha = 0.05$), tingkat \textit{Improved \%} mencapai 100\% pada dua nilai $\alpha$ (0.05 dan 0.01), artinya seluruh graf mengalami perbaikan dari kondisi awalnya. Hal ini mengindikasikan bahwa topologi VHS yang lebih sederhana (7 node) lebih mudah diperbaiki secara relatif, meskipun ambang keseimbangan tetap sulit dicapai karena tidak ada satu pun graf yang seimbang pada kondisi awal (\textit{Initial Balanced} = 0\%). \textit{Fitness Std Dev} yang sangat rendah (0.0507 pada $\alpha = 0.05$) juga menunjukkan konsistensi hasil yang tinggi pada dataset ini.

Analisis distribusi per tipe agen pada dataset abstrak ($\alpha = 0.05$) disajikan pada Tabel \ref{tab:per_agent} dan Gambar \ref{fig:per_agent_plot}.

\begin{table}[htbp]
\centering
\caption{Perbandingan Hasil per Profil Agen pada Dataset Abstrak ($n=30$, $\alpha=0.05$, Kondisi C)}
\label{tab:per_agent}
\begin{tabular}{lccc}
\hline
\textbf{Profil Agen} & \textbf{Balanced \%} & \textbf{Improved \%} & \textbf{Median Gen} \\
\hline
\textit{Aggressive} & 63.33\% & 83.33\% & 3.0 \\
\textit{Passive}    & 60.00\% & 83.33\% & 6.5 \\
\textit{Random}     & 56.67\% & 73.33\% & 6.5 \\
\hline
\end{tabular}
\end{table}

\begin{figure}[htbp]
    \centering
    \includegraphics[width=0.75\textwidth]{tugas-akhir/assets/images/plot_per_agent.png}
    \caption{Perbandingan Distribusi Hasil per Profil Agen ($\alpha=0.05$)}
    \label{fig:per_agent_plot}
\end{figure}

Agen \textit{aggressive} secara konsisten menghasilkan performa terbaik dengan \textit{Balanced \%} 63.33\% dan konvergensi tercepat (median generasi 3.0). Hal ini disebabkan oleh interaksi agen agresif yang lebih sering memicu \textit{Converter} manual, menghasilkan aliran sumber daya yang lebih aktif dan sinyal fitness yang lebih kaya bagi algoritma evolusioner. Agen \textit{passive} dan \textit{random} menunjukkan \textit{Improved \%} yang sama dengan \textit{aggressive} (83.33\%) namun konvergensi lebih lambat (6.5 generasi), mencerminkan tantangan ekstra dari perilaku pemain yang lebih konservatif dan stokastik.

\subsection{Perbandingan dengan Penelitian Dasar}
\label{subsec:paper_comparison}

Tabel \ref{tab:paper_comparison} membandingkan hasil penelitian ini dengan hasil paper GEEvo asli \citep{Rupp_2024} sebagai referensi.

\begin{table}[htbp]
\centering
\caption{Perbandingan Hasil Penelitian dengan Paper GEEvo \citep{Rupp_2024}}
\label{tab:paper_comparison}
\begin{tabular}{lcccc}
\hline
\textbf{Sumber} & \textbf{Balanced \% ($\alpha$=0.05)} & \textbf{Balanced \% ($\alpha$=0.01)} & \textbf{Median Gen} & \textbf{Exec Time (s)} \\
\hline
Paper GEEvo (n=194)    & 93.3\%  & 83.0\% & 1.0   & 18.4   \\
TA — Abstrak (n=30)    & 66.67\% & 30.0\% & 2.5   & 1.191  \\
TA — MMORPG (n=10)     & 50.0\%  & 0.0\%  & 251.0 & 3.609  \\
TA — VHS/Horror (n=10) & 30.0\%  & 20.0\% & 500.0 & 18.004 \\
\hline
\end{tabular}
\end{table}

Perbedaan hasil dengan paper asli bersifat by design dan dapat dijelaskan oleh tiga faktor metodologis:
\begin{itemize}
    \item \textbf{Ukuran dataset}: Paper menggunakan $n = 194$ graf, sedangkan penelitian ini menggunakan $n = 30$ (abstrak) dan $n = 10$ (GDD), sehingga estimasi statistik kurang stabil dan rentan terhadap variasi sampel.
    \item \textbf{Kompleksitas topologi GDD}: Dataset MMORPG dan VHS memiliki topologi yang ditentukan dari konten dokumen nyata, yang tidak selalu memiliki karakteristik optimal untuk penyeimbangan cepat.
    \item \textbf{Penambahan simulasi agen}: Penggunaan \texttt{AgentSimulator} menambahkan lapis stokastisitas tambahan yang tidak ada pada paper asli, sehingga mempersulit konvergensi namun meningkatkan representasi pemain nyata.
\end{itemize}

Perbandingan ini bersifat \textit{referensial} — fokus utama evaluasi penelitian ini adalah perbandingan antar Kondisi A, B, dan C, bukan replikasi eksak paper asli. Tabel \ref{tab:paper_comparison} menunjukkan bahwa pada dataset abstrak ($n=30$), sistem TA mencapai \textit{Balanced \%} sebesar 66.67\% dengan waktu eksekusi rata-rata hanya 1.191 detik per graf, jauh lebih cepat dari paper (18.4 detik) meskipun dengan $\alpha$ yang sama.

    \chapter{Kesimpulan dan Saran}
    \section{Kesimpulan}

        Berdasarkan hasil implementasi dan pengujian yang telah dilakukan, penelitian ini menarik beberapa kesimpulan sebagai jawaban atas rumusan masalah yang diajukan:

        \begin{enumerate}
            \item Model sistem \textit{reward} berbasis graf berarah $G = (V, E, W)$ berhasil diimplementasikan untuk merepresentasikan aliran sumber daya dalam ekonomi gim. Lima tipe node — \textit{Source}, \textit{Pool}, \textit{Converter}, \textit{RandomGate}, dan \textit{Drain} — bersama-sama membentuk ekosistem aliran sumber daya yang dapat dimodelkan, disimulasikan, dan dioptimasi secara komputasional. Validasi melalui simulasi deret waktu menunjukkan bahwa model mampu merepresentasikan kondisi tidak seimbang (akumulasi atau penipisan sumber daya) maupun kondisi seimbang (aliran stasioner dinamis) sesuai dengan konfigurasi bobot \textit{edge} yang ditetapkan.

            \item Simulasi berbasis agen dengan tiga profil perilaku (\textit{aggressive}, \textit{passive}, \textit{random}) berhasil dikembangkan untuk merepresentasikan variabilitas perilaku pemain. Evaluasi komparatif pada $n = 194$ graf abstrak menunjukkan bahwa penambahan simulasi agen (Kondisi B vs Kondisi A) secara konsisten meningkatkan kualitas pencarian algoritma evolusioner: \textit{Improved \%} naik sebesar $+3.09\%$ dan \textit{Fitness Std Dev} turun sebesar 0.049, mengindikasikan bahwa agen membantu algoritma menemukan solusi yang lebih baik dengan variansi yang lebih rendah. Analisis per profil agen menunjukkan bahwa agen \textit{aggressive} memberikan performa terbaik (\textit{Balanced \%} = 63.33\%, Median Generasi = 3.0) karena interaksinya yang aktif menghasilkan sinyal fitness yang lebih informatif bagi algoritma.

            \item Integrasi antara simulasi berbasis agen dan fungsi fitness obj4 (\textit{tolerance-threshold}) berhasil dikembangkan sebagai kontribusi utama penelitian ini. Kondisi C (agen + obj4) menunjukkan keunggulan signifikan pada $\alpha = 0.05$: \textit{Balanced \%} mencapai 54.64\% (lebih tinggi $+8.25\%$ dari Kondisi B), \textit{Improved \%} mencapai 87.11\%, dan \textit{Median Generation} hanya 8 generasi dibandingkan 500 generasi pada Kondisi A dan B. Konvergensi yang jauh lebih cepat pada Kondisi C mengkonfirmasi bahwa obj4 memberikan gradien sinyal yang lebih efektif bagi algoritma evolusioner. Validasi pada dataset GDD nyata (MMORPG dan VHS) menunjukkan bahwa pendekatan ini dapat diterapkan pada ekonomi gim nyata yang diekstrak secara otomatis dari \textit{Game Design Document}: seluruh graf VHS ($n=10$) mengalami perbaikan (Improved \% = 100\%) meskipun ambang keseimbangan ketat lebih sulit dicapai pada topologi yang lebih kompleks.
        \end{enumerate}

    \section{Saran}

        Berdasarkan pengalaman dalam penelitian ini, beberapa saran diajukan untuk pengembangan lebih lanjut:

        \begin{enumerate}
            \item Penelitian lanjutan dapat mempertimbangkan penggunaan dataset yang lebih besar ($n \geq 100$) untuk dataset GDD agar estimasi statistik lebih representatif dan stabil. Hal ini memerlukan ketersediaan GDD yang lebih banyak atau augmentasi GDD yang ada melalui variasi konfigurasi topologi.

            \item Fungsi fitness obj4 terbukti efektif pada $\alpha = 0.05$ namun konvergen identik dengan obj3 pada $\alpha$ yang lebih ketat ($\alpha \leq 0.01$). Penelitian lanjutan dapat mengeksplorasi formula fitness adaptif yang menyesuaikan mekanisme \textit{tolerance-threshold} secara dinamis berdasarkan kemajuan konvergensi, sehingga manfaat obj4 dapat dipertahankan pada kondisi $\alpha$ yang lebih ketat.

            \item Modul \texttt{GDDExtractor} saat ini menggunakan pendekatan pencocokan kata kunci statis. Penggunaan model bahasa besar (\textit{large language model}) untuk pemahaman semantik dokumen GDD secara kontekstual dapat meningkatkan akurasi pemetaan node, terutama untuk GDD yang menggunakan terminologi tidak konvensional atau dalam bahasa selain Inggris.

            \item Penelitian ini dapat diperluas untuk mengoptimasi topologi graf secara bersamaan dengan parameter bobot (\textit{co-optimization}) menggunakan algoritma evolusioner multi-objektif, sehingga tidak hanya bobot \textit{edge} tetapi juga struktur koneksi antar node dapat diadaptasi terhadap profil perilaku pemain.
        \end{enumerate}

    \chapter{Lampiran}\label{ch:appendix} % Bisa dikasih `comment` kalau tidak dipakai
    \input{tugas-akhir/contents/007_bibliography}
    \chapter*{Halaman Persembahan}
Penyusunan Tugas Akhir ini dapat terselesaikan tidak terlepas dari berbagai pihak yang telah memberikan dukungan, motivasi, dan bantuan secara langsung maupun tidak langsung. Oleh karena itu, penulis ingin menyampaikan ucapan terima kasih kepada: 

%
% \begin{enumerate}
%     \item Kedua orang tua ....
%     \item Diri sendiri selaku penulis, atas perjuangan dan dedikasi hingga dapat menyelesaikan Tugas Akhir ini tepat waktu.
%     \item Mr B
%     \item Teman-teman pejuang TA : .... yang telah menjadi rekan diskusi dan saling membantu selama ....
%     \item Kawan-kawan Jogging Keras \& KOMPAS yang telah memotivasi pengerjaan Tugas Akhir ini sekaligus teman-teman Penghuni Lab: ... dan semua yang tidak dapat disebutkan satu per satu yang telah menjadi tempat berbagi suka dan duka selama masa perkuliahan.   
%     % \item Seseorang yang selalu aku kagumi, yang mendorongku untuk terus menjadi lebih baik dari hari kemarin.
%     \item Teman-teman Departemen Matematika ITS angkatan 2022 AXXIOM yang telah menemani masa-masa perkuliahan penulis.
%     \item Siapapun
%     \item Cup of Joe, The 1975, ...
%     \item Pihak lain yang berjasa dalam membantu penyelesaian \ta~yang belum dapat disebutkan.
% \end{enumerate}
%
Penulis juga mengharapkan kritik dan saran yang membangun dari berbagai pihak untuk penyempurnaan isi Tugas Akhir ini. Akhir kata, semoga Tugas Akhir ini bermanfaat bagi semua pihak yang bersangkutan.
% \input{tugas-akhir/contents/halaman_kosong}
\clearpage
\chapter*{Biodata Penulis}
    % Foto pribadi 
    \begin{wrapfigure}{l}{0.3\textwidth}
        
\includegraphics[width = \linewidth]{tugas-akhir/assets/images/gojo.png} % Ganti fotonya sesuai tempat direktorinya 
    \end{wrapfigure}
    Nama lengkap penulis Alridho Trisan Satriawan, lahir di Kabupaten Pati pada 11 Juli 2002. Sebelum menempuh bangku perkuliahan, penulis memulai pendidikan formal mulai dari TK RA Masyitoh (2007-2009), SDN Tegalharjo 01 (2009-2015), SMPN 3 Pati (2015-2018), SMAN 1 Juwana (2018-2021). Pada tahun 2022, penulis meneruskan pendidikan ke jenjang strata-1 di Departemen Matematika Institut Teknologi Sepuluh Nopember (ITS). Pada masa perkuliahan, penulis aktif dalam kegiatan ekstra kampus antara lain, anggota luar biasa UKM GIM ITB sebagai gaem designer (2023-2024), anggota komunitas Level UP sebagai divisi web \& copywriting (2024-2025), anggota komunitas DigistarClub dari Telkom sebagai divisi creative (2024-2025). Penulis juga aktif dalam bidang lingkungan dan teknologi dengan mengikuti program ambassadorship AIVolution yang diselenggarakan oleh kedutaan besar Amerika Serikat. Di bidang keahlian, penulis memiliki keterampilan pada analisis data, pemrograman (Java, Python), desain antarmuka, pengembangan gim, serta keamanan siber. Sejumlah sertifikasi dari IBM, Cisco, dan Dicoding turut memperkuat kompetensi penulis dalam teknologi dan inovasi digital. Adapun informasi lebih lanjut dapat menghubungi penulis melalui email alridho.tristan@gmail.com.
\end{document}